% concatenated from MN-26.01.26.tex
 \documentclass[11pt]{amsart}
\usepackage{amsmath,amsfonts,amssymb,upgreek, amscd, mathtools}
\setlength{\unitlength}{0.4in}
\renewcommand{\baselinestretch}{1.2}
\setlength{\topmargin}{-9mm}
%\setlength{\topmargin}{-20mm}
\setlength{\textheight}{9.2in}
\setlength{\textwidth}{6.6in}
\setlength{\oddsidemargin}{-0.1in}
\setlength{\evensidemargin}{-0.1in}

\setlength{\abovedisplayshortskip}{3mm}
\setlength{\belowdisplayshortskip}{3mm}
%\renewcommand{\baselinestretch}{1.5}
%\vskip.1in


\usepackage[breaklinks]{hyperref}

\theoremstyle{thrm}
\usepackage{tikz}
\usepackage{enumerate}
\usetikzlibrary{calc,decorations.markings}
\usepackage{tikz-cd}
%\usetikzlibrary{cd}
\usepackage{graphicx}


\theoremstyle{plain}
\newtheorem{no}{Notation}[section]
\newtheorem{Prob}{Problem}
\newtheorem{thm}{Theorem}[section]
\newtheorem{lemma}[thm]{Lemma}
\newtheorem{prop}[thm]{Proposition}
\newtheorem*{Algo}{Algorithm}
\newtheorem{cor}[thm]{Corollary}
\newtheorem{con}[thm]{Conjecture}
\newtheorem{question}[thm]{Question}
\newtheorem{defn}[thm]{Definition}
\newtheorem{example}[thm]{Example}
\newtheorem*{ack}{Acknowledgement}
\newtheorem{remark}[thm]{Remark}


\theoremstyle{definition}
\newcommand{\Ad}{\operatorname{Ad} }
\newcommand{\e}{\operatorname{e} }
\newcommand{\ad}{\operatorname{ad} }
\newcommand{\Der}{\operatorname{Der} }
\newcommand{\A}{\operatorname{A} }
\newcommand{\RBO}{\operatorname{RBO} }
\newcommand{\Sub}{\operatorname{Sub} }
\newcommand{\Exp}{\operatorname{Exp} }
\newcommand{\B}{\operatorname{B} }
\newcommand{\C}{\operatorname{C} }
\newcommand{\D}{\operatorname{D} }
\newcommand{\E}{\operatorname{E} }
\newcommand{\F}{\operatorname{F} }
\newcommand{\K}{\operatorname{K} }
\newcommand{\Ho}{\operatorname{H} }
\newcommand{\M}{\operatorname{M} }
\newcommand{\N}{\operatorname{N} }
\newcommand{\RH}{\operatorname{RH} }
\renewcommand{\O}{\operatorname{O} }
\newcommand{\Z}{\operatorname{Z} }
\newcommand{\Cl}{\operatorname{Cl} }
\newcommand{\cl}{\operatorname{cl} }
\newcommand{\im}{\operatorname{im} }
\newcommand{\GF}{\operatorname{GF}}
\newcommand{\Fix}{\operatorname{Fix}}
\newcommand{\het}{\operatorname{ht}}
\newcommand{\Hol}{\operatorname{Hol}}
\newcommand{\Id}{\operatorname{id}}
\newcommand{\Sh}{\operatorname{Sh}}
\renewcommand{\Pr}{\operatorname{Pr} }
\newcommand{\I}{\operatorname{I} }
\renewcommand{\det}{\operatorname{det} }
\newcommand{\gen}[1]{\left < #1 \right >}
\newcommand{\Aut}{\operatorname{Aut} }
\newcommand{\Autb}{\operatorname{Autb} }
\newcommand{\Tra}{\operatorname{Tra} }
\newcommand{\Res}{\operatorname{Res} }
\newcommand{\Inf}{\operatorname{Inf} }
\newcommand{\Deg}{\operatorname{Deg} }
\newcommand{\End}{\operatorname{End} }
\newcommand{\Ext}{\operatorname{Ext} }
\newcommand{\STExt}{\operatorname{STExt} }
\newcommand{\CExt}{\operatorname{CExt} }
\newcommand{\Gal}{\operatorname{Gal} }
\newcommand{\Hom}{\operatorname{Hom} }
\newcommand{\Fun}{\operatorname{Fun} }
\newcommand{\Map}{\operatorname{Map} }
\newcommand{\Inn}{\operatorname{Inn} }
\newcommand{\Irr}{\operatorname{Irr} }
\newcommand{\Ker}{\operatorname{ker} }
\newcommand{\Out}{\operatorname{Out} }
\newcommand{\Ann}{\operatorname{Ann} }
\newcommand{\Soc}{\operatorname{Soc} }
\newcommand{\IM}{\operatorname{im} }
\newcommand{\Lin}{\operatorname{Lin} }
\newcommand{\Mat}{\operatorname{Mat} }
\newcommand{\Syl}{\operatorname{Syl} }
\newcommand{\Autcent}{\operatorname{Autcent} }
\newcommand{\into}{\rightarrowtail}
\newcommand{\onto}{\twoheadrightarrow}
\newcommand{\normal}{\vartriangleleft}
\newcommand{\normaleq}{\trianglelefteq}
\newcommand{\annd}{\quad \text{ and } \quad}
\newcommand{\inn}{\text{\textrm{\textup{ in }}}}
\newcommand{\onn}{\text{\textrm{\textup{ on }}}}
\newcommand{\orr}{\quad \text{ or } \quad}
\newcommand{\br}[1]{\overline{#1} }
\newcommand{\htt}[1]{\widehat{#1} }
\newcommand{\tld}[1]{\widetilde{#1} }
\numberwithin{equation}{section}
\newcommand{\op}{\operatorname{op} }
\newcommand{\Diag}{\operatorname{Diag} }

\begin{document}

\title{Skew brace extensions, second cohomology and  complements}
	

\author{Nishant Rathee}
\address{ Department of Mathematics, National Iinstitute of Technology, Tiruchirappalli-620015, India}
\email{nishant@nitt.edu, monurathee2@gmail.com}		

\author{Manoj K. Yadav}
\address{Harish-Chandra Research Institute, A CI of Homi Bhabha National Institute, Chhatnag Road, Jhunsi, Prayagraj-211 019, India}
\email{myadav@hri.res.in}

\subjclass[2010]{16T25, 81R50, 20J06}
\keywords{Skew left brace,  ideal, annihilator, socle, extension, second cohomology, complement}

\begin{abstract}
We study extensions and the second cohomology groups of skew left braces via the associated natural semi-direct product groups.  Let $0 \to I \to E \to H \to 0$ be a skew brace extension and $\Lambda_H$ denote  the natural semi-direct product associated with the skew left brace $H$. We establish a group homomorphism from $\Ho_{Sb}^2(H, I)$, the second cohomology group of braces, into $\Ho_{Gp}^2(\Lambda_H, I \times I)$, the second cohomology group of groups, which turns out to be an embedding when $I \le \Soc(E)$, the socle of $E$. In particular, the Schur multiplier of a skew left brace  $H$ embeds into the Schur multiplier of the group $\Lambda_H$. Analog of the Schur-Zassenhaus theorem is established for skew left braces in several specific cases. We introduce a concept called minimal  extensions (which stay at the extreme end of split extensions)  of skew left braces and derive many fundamental results.  Several reduction results for split extensions of finite skew left braces by abelian groups (viewed as  trivial left braces)  are obtained. 

\end{abstract}


\maketitle
	
\section{Introduction}

Cohomology of left braces and linear cycle sets with trivial action on  an abelian group was introduced  by Labed and Vendramin  \cite{LV16} .  This was extended to the case when the action is non-trivial by Bardakov and Singh \cite{BS21}. The concept of the (first and second) cohomology of skew left braces,  originated from the article \cite{DB18}  by Bachiller, was formalised by the authors of the present article in \cite{RY24}.  Connection with extension theory was developed. The aim of this article is to further advance this knowledge by exploiting the connection of skew left brace with the natural semi-direct product of the multiplicative group of the skew left brace by that of its additive group. Several foundational results on minimal extensions of skew left braces and Frattini sub-skew braces are produced. Some reduction techniques for an exact sequence of skew left braces  by abelian $p$-groups to split are obtained.
	

A skew left brace is a set \(A\) equipped with two group structures, \((A, +)\) and \((A, \circ)\) satisfying the compatibility condition
\[
a\circ (b  +  c) = (a \circ b) - a + (a \circ c),\]
for all \( a,b,c \in A\), where \(- a\) denotes the inverse of the element $a$ in the group \((A, +)\).  A skew left brace $A$ is denoted by $(A, +, \circ)$, mostly when the binary operations are emphasized.  The groups  \((A, +)\) and \((A, \circ)\), respectively, termed as the additive group and multiplicative group of the skew left brace $A$.


Given a skew left brace $(A, +, \circ)$, there exists an action of $(A, \circ)$ on $(A, +)$ by automorphisms given by
\[
\lambda: (A, \circ) \to \operatorname{Aut} (A, +), \quad \lambda_a(b) = -a + (a \circ b),
\]
which is known as the lambda map associated with skew left brace $A$. This action allows one to construct the associated semi-direct product group 
\[
\Lambda_A = (A, +) \rtimes_{\lambda} (A, \circ).
\]
Let $\Ho^2_{Sb}(H, I)$ denote the second brace cohomology group of a skew left brace $H$ with coefficients in an abelian group $I$ and  $\Ho^2_{Gp}(\Lambda_H, I \times I)$  denote the second group cohomology group of the associated group $\Lambda_H$ with coefficients in the  abelian group $I \times I$.
Our first result is
\begin{thm}
Let a skew left brace $H$ act on an abelian group via an action $(\nu, \mu, \sigma)$. There exists a group homomorphism  $\Psi : \Ho^2_{Sb}(H, I) \to \Ho^2_{Gp}(\Lambda_H, I \times I)$. Moreover, if   $\mu$ is trivial and $\nu^{-1}_h = \sigma_h$ for all $h \in H$, then $\Psi$ is injective.
\end{thm}

As a quick application of this theorem we get \cite[Theorem 2.2]{LV24} which states that the exponent of the Schur multiplier of a finite skew left brace $H$ divides $|H|^2$. 


Our next result is about the connection between the  second cohomology group of a skew left brace $H$ and that of its opposite skew left brace $H^{\op}$.

\begin{thm}
Let $H$ be a skew left brace acting on a  trivial brace $I$ via a good triplet of actions $(\nu, \mu, \sigma)$. Then $\Ho^2_{Sb}(H, I) \cong \Ho^2_{Sb}(H^{\op}, I)$. 
\end{thm}



An ideal $I$ of a finite skew left brace $H$  is said to be a Hall ideal if the order of $I$ is coprime to the order of the quotient skew left brace $H/I$. We prove a brace analogous of the group theory result \cite[9.1.2]{BEJP24} in several specific cases (Theorem \ref{main7}), for example,  when (i) $I \le \Ker(\lambda)$, (ii)  $E$ is supersoluble, (iii) $(E, +)$ is a nilpotent group. We present various  interesting results on Frattini  sub-skew braces, minimal and split extensions of skew left braces. For example, we prove that the Frattini sub-skew  brace of the direct sum of two finite skew left braces $E_1$ and $E_2$ is the direct sum of the Frattini sub-skew braces of $E_1$ and $E_2$.

In the last section we prove a reduction result for split  brace extensions.
\begin{thm}\label{thm1}
	 Let $E$ be a left brace  admitting an ideal  $I$ such that $I = I_1 \times \cdots \times I_r$, where  $I_j$ is an ideal of $E$ for each $1 \leq j \leq r$. Let $I_k^* = \prod_{j \neq k} I_j$. Then the following  hold true:
\begin{itemize}
	\item[(i)] $E$ splits over $I$ if and only if $E/I_k^*$ splits over $I/I_k^*$ for all $k = 1, \ldots, r$.
	\item[(ii)]  $E$ splits over $I$ if and only if, for any $k \neq m$, the quotients $E/I_k$ and $E/I_m$ split over $I/I_k$ and $I/I_m$, respectively.
\end{itemize}
\end{thm}

 We conclude this section by setting some notations. Note that additive and multiplicative identities of any skew left brace coincide, which we denote by $0$. A skew left brace having cardinality $1$ is also mostly denoted by $0$, instead of $\{0\}$.  Inverse of  an element $a \in A$, when considered in $(A, +)$  (respectively in $(A, \circ)$)  is denoted by $-a$  (respectively by $a^{-1}$). The group commutator of two elements $a$ and $b$ of a skew left brace $A$ in $(A, +)$ is denoted  by $[a, b]^{+} := -a - b + a + b$, and in $(A, \circ)$, by $[a, b]^{\circ} := a^{-1} \circ b^{-1} \circ  a \circ b $. The commutator  $[a, b]^{+}$ is also denoted by $\gamma_{+}(a, b)$. The commutators of larger weights are defined iteratively.  We write $B \le A$ to denote that $B$ is a sub-skew brace of the skew  brace $A$, and $B \trianglelefteq A$ to denote that $B$ is an ideal  of  $A$.   The quotient skew brace  $A/I$ of $A$ by $A$ will be denoted by $\bar{A}$. Isomorphism theorems for skew braces were established in \cite{BEP24}, which we'll use without further reference. The category of all skew left braces will be denoted by $\mathcal{SB}$.


\section{Preliminaries}

A {\it skew left brace} is a set \(A\) equipped with two group structures, \((A, +)\) and \((A, \circ)\) satisfying the compatibility condition
\[a\circ (b  +  c) = (a \circ b) - a + (a \circ c),\]
for all \( a,b,c \in A\), where \(- a\) denotes the inverse of the element $a$ in the group \((A, +)\).  A skew left brace $A$ is denoted by $(A, +, \circ)$, mostly when the binary operations are emphasized.  The groups  \((A, +)\) and \((A, \circ)\), respectively, termed as the additive group and multiplicative group of the skew left brace $A$.  A {\it right skew brace} is defined analogously, demanding the  `skew-distributivity' from the right. A skew left brace which is also a skew right brace is a called a {\it two-sided skew brace}. A skew left brace is said to be a {\it left brace} if $(A, +)$ is an abelian group.  We say that a skew left brace is {\it trivial} if `$+$' and `$\circ$' coincide on $A$. So a trivial skew left brace $(A, +, \circ)$ is just the group $(A, +)$ and a trivial left  brace $(A, +, \circ)$ is just the abelian group $(A, +)$. As we always work with skew left braces, for  simplicity, we shall  suppress the use of the word `left', and just call skew brace. Likewise, a left braces will be called a brace.  A `trivial brace' and a `trivial skew brace of abelian type' mean the  same thing.


For a skew  brace $(A, +, \circ)$, it was  proved in \cite{GV17} that the map
$$
\lambda  :  (A, \circ) \to \Aut \, (A, +),~~a \mapsto \lambda_a
$$
is a group homomorphism, where $\Aut \, (A, +)$ denotes the automorphism  group of $(A, +)$ and $\lambda_a$ is given by $\lambda_a(b) = - a + (a \circ b)$ for all $a, b \in A$.  
For two elements $a, b \in A$,  $-b + a + b$ denotes the conjugate of $a$ by $b$ in $(A, +)$ and    $ b^{-1} \circ a \circ b$ denotes  the conjugate of $a$ by $b$ in $(A, \circ)$. 

A subset $B$ of a skew  brace $A = (A, +, \circ)$ is said to be a {\it sub-skew brace} of  $A$ if $B$ is a subgroup of both the group structures $(A, +)$ and $(A, \circ)$. A subset $I$ of a skew  brace $(A, +, \circ)$ is said to be a \emph{ left ideal } of $(A, +, \circ)$ if $I$ is a  subgroup of $(A, +)$ and $\lambda_a(c) \in I$ for all $a \in A$ and $c \in I$. It turns out that a left ideal $I$ of $(A, +, \circ)$ is a subgroup of $(A, \circ)$. A left ideal $I$ of $A$ is said to be  a \emph{strong left ideal} if $I$ is normal in $(A, +)$. A left ideal $I$ is said to be an \emph{ideal} of  $(A, +, \circ)$ if $I$ is normal in both $(A, +)$ and  $(A, \circ)$.  Note that $\Ker(\lambda)$ is a left  ideal of $(A, +, \circ)$.


Let $A := (A, +, \circ)$ be a skew  brace. We define {\it socle} of $A$ to be the set 
$$\Soc(A) := \Ker(\lambda) \cap \Z(A, +)$$
 and {\it annihilator} of $A$ to be the set 
 $$\Ann(A) :=  \Ker(\lambda) \cap \Z(A, +) \cap \Z(A, \circ).$$ 
 It turns out that both $\Soc(A)$ and $\Ann(A)$ are ideals of $A$.



The following result is well known.
\begin{prop}
	Let $A$ be a brace and let $S$ and $I$ be, respectively, a sub-skew brace and an ideal of $A$. Then $S \circ I = S + I$ is a sub-skew brace of $A$. Moreover, if $S$ is  a left ideal (an ideal), then $S + I$ is also  a left ideal (an ideal).
\end{prop}

	Let $A$ be a brace and $B$  a sub-skew brace of $A$. We say that $B$ is {\it supplemented} in $E$ if  there exists a sub-skew brace $C$ such that
$$A = B \circ C =  B + C.$$
Such a $C$ is called a supplement of $B$ in $E$.
The sub-skew brace  $B$ is said to be  \emph{complemented} in $E$ if $B$ is supplemented by  $C$ and 
$B \cap C = \{0\}$. In this case we say that $C$ is a complement of $B$ in $E$.
Note that if $A =  B + C$, where $B$  is a skew brace of $E$ and $C$ is an ideal of $E$, then $A = C + B = B \circ C = C \circ B$. Indeed, $a = b + c = b \circ \lambda^{-1}_b(c)$ for all $a, b, c \in A$. Thus, an ideal $I$ of $A$ is complemented by a skew brace $B$ in $A$ if and only if $A = I + B$ and $I \cap B =  \{0\}$.

Let $H$ and $I$ be two skew braces. By a  \emph{ brace extension} of $H$ by $I$, we mean a skew brace $E$ with an exact sequence
\[
0 \longrightarrow I \xrightarrow{i} E \xrightarrow{\pi} H \longrightarrow 0,
\]
where $i$ and $\pi$ are injective and surjective skew brace homomorphisms, respectively.  We'll mostly view $I$  as an ideal $E$ via $i$. When we do not wish to emphasize the use of the maps $i$ and $\pi$,  we say that  $E$ is a brace extension of $H$ by $I$ without referring to the exact sequence.  A set map $s: H \to E$ is said to be an {\it st-section} if $s(0) = 0$ and $s \pi = \Id_H$.   An st-section $s$ is said to be a {section} if it is a skew brace homomorphism. The preceding extension $E$ is said to be a {\it split extension} if there exists a  section $s: H \to E$. 


Let us start with a brace extension of a skew brace $H$ by a trivial brace $I$
\[
0 \longrightarrow I \xrightarrow{i} E \xrightarrow{\pi} H \longrightarrow 0
\]
with  an st-section $s : H \to E$.  Then the multiplicative group $(H, \circ)$ acts on the additive group $(I, +)$ from the left via an action $\nu$ defined by 
$$\nu(h)(x) = \lambda_{s(h)}(x)$$
 for all $h \in H$ and $x \in I$. The groups $(H, +)$ and $(H, \circ)$ act  on $(I, +)$ from the right via actions $\mu$ and $\sigma$, defined, respectively,  
 $$\mu(h)(x) = -s(h) + x + s(h)$$
  and
  $$\sigma(h)(x) =  s(h)^{-1} \circ x\circ s(h)$$
  for all $h \in H$ and $x \in I$. The images of $h \in H$ in $\operatorname{Aut}(I, +)$, the automorphism group of the additive group $(I, +)$, under the actions $\nu$, $\mu$, and $\sigma$ are denoted by $\nu_h$, $\mu_h$, and $\sigma_h$, respectively. The triple of actions $(\nu, \mu, \sigma)$ satisfy the following compatibility conditions:
\[
\mu_{\lambda_{h_1}(h_3)}(y) = \nu_{h_1} ( \mu_{h_3} ( \nu_{h_1}^{-1}(y) )),
\]
\[
\nu_{h_2 + h_3} \left( \sigma_{h_2 + h_3}(y) \right) + \mu_{h_3}(y) = \mu_{h_3} \left( \nu_{h_2} \left( \sigma_{h_2}(y) \right) \right) + \nu_{h_3} \left( \sigma_{h_3}(y) \right),
\]
for all $h_1, h_2, h_3 \in H$ and $y \in I$. In general, we say that a skew brace $H$ {\it acts on an abelian group} $I$, viewed as a trivial brace, via a good triplet of actions $(\nu, \mu, \sigma)$ if $\mu:  (H, +) \to \Aut(I)$ and $\sigma: (H, \circ) \to  \Aut(I)$  are group anti-homomorphisms and $\nu: (H, \circ) \to \Aut(I)$ is a group homomorphism and satisfy the preceding two compatibility conditions, where $\nu_h$, $\mu_h$, and $\sigma_h$, respectively, are the images of $h \in H$ under the actions $\nu$, $\mu$, and $\sigma$.


As explained above $\lambda$ is an action of $(A, \circ)$ on $(A, +)$. This action allows us to construct the associated semi-direct product group 
\[
\Lambda_A = (A, +) \rtimes_{\lambda} (A, \circ),
\]
in which product of two elements $(x_1, y_1)$ and $(x_2, y_2)$ is given by $(x_1 + \lambda_{y_1}(x_2), y_1 \circ y_2)$.
It turns out that the association $A \mapsto \Lambda_A$ is a functor from the category of skew braces to the category of groups. Therefore a skew brace homomorphism between two skew  braces $\phi: E \rightarrow H$ defines a group homomorphism between the corresponding groups $\Lambda_E$ and $\Lambda_H$, given by
\[
(\phi \times \phi) : \Lambda_E \rightarrow \Lambda_H, \quad (\phi \times \phi)(x, y) = (\phi(x), \phi(y)) \]
for all $x, y \in E$. Thus, every brace extension 
\begin{equation}\label{sb-ext}
0 \longrightarrow I \xrightarrow{i} E \xrightarrow{\pi} H \longrightarrow 0,
\end{equation}
induces the following group extension:
\begin{equation}\label{g-ext}
0 \longrightarrow \Lambda_I  \xrightarrow{i \times i} \Lambda_E \xrightarrow{\pi \times \pi} \Lambda_H \longrightarrow 0.
\end{equation}

The brace extensions $E$ and $E'$  of a skew brace $H$ by a skew brace $I$ are said to be \emph{equivalent} if there exists a brace homomorphism $\phi : E \to E'$ such that the following diagram commutes:
$$\begin{CD}
 0 @>i>> I @>>> E @>{{\pi} }>> H  @>>> 0\\
 &&  @V{\text{Id}}VV @V{\phi}VV @ VV{ \text{Id}}V \\
 0 @>i'>> I @>>> E^\prime @>{{\pi^\prime} }>> H @>>> 0.
\end{CD}$$
The set of all equivalence classes of extensions of $H$ by $I$ is denoted by $\Ext(H, I)$. When we want to emphasize the use of an action $(\nu, \mu, \sigma)$ of $H$ on $I$, then we denote $\Ext(H, I)$ by $\Ext_{(\nu, \mu, \sigma)}(H, I)$.




\begin{prop}
Let \( E \) be a skew brace extension of \( H \) by \( I \). Then the following hold true:
\begin{itemize}
	\item[(i)] If \( E \) is a split extension, then so is  \( \Lambda_E \).
	\item[(ii)] If \( E \) is an  extension of \( H \) by  a trivial brace \( I \), then the extension \( \Lambda_E \) is an abelian extension of \( \Lambda_H \) by \( I \times I \). Furthermore, if  $I \leq \Ann(E)$, then the induced extension of $\Lambda$-groups is a central extension as $\Lambda_{\Ann(E)} \leq \Z(\Lambda_E)$ (see \cite[Proposition 3.4]{RSU24} for details).
	\item[(iii)]  If \( E \) and \( E' \) are equivalent skew brace extensions of \( H \) by \( I \), then \( \Lambda_{E} \) and \( \Lambda_{E'} \) are also equivalent group extensions of \( \Lambda_H \) by \( \Lambda_I \).
\end{itemize}
\end{prop}

So it turns out that  the covariant functor $\Lambda$ from the category of skew braces to the category of groups is exact. 



We conclude  this section with a couple of definitions, which will be used in the last two sections. Let $H$ be a skew  brace and $I$  an abelian group, on which $H$ acts via a good triplet of actions $(\nu, \mu, \sigma)$. In such a case, we'll say that $I$ is an {\it $H$-module}.   A subgroup \(M \leq I\) is said to be  a {\it submodule} if it is invariant under the actions \((\nu, \mu, \sigma)\). Furthermore, an \(H\)-module is called \emph{simple} if it has no proper non-zero submodules.




%%%%%%%%%%%%%%%%%%%%%%%%%%%%%%%%%%%%%%%%%%%

\section{From skew braces to groups : extensions and cohomology}


Let $H$ be a skew brace and $I$  a trivial brace such that $H$ acts on $I$ via a good triplet of actions $(\nu, \mu, \sigma)$. We start by defining an action of $\Lambda_H$ on $I \times I$. Let $E$ be any brace extension of $H$ by $I$. Then, as explained above,  $\Lambda_E$ turns out to be an extension of $\Lambda_H$ by $I \times I$. We know, again as explained above, that the actions $(\nu, \mu, \sigma)$ can be completely determined via the extension $E$. Using that characterization, for  $x \in I$ and $h \in H$, we get
\begin{align}\label{action1}
	\lambda_x(s(h))-s(h) &= -x+ x \circ s(h)-s(h) \nonumber \\
	&=-x+ s(h) \circ \sigma_h(x)-s(h) \nonumber\\
	&=-x+ s(h)+\nu_h(\sigma_h(x))-s(h) \nonumber\\
	&=-x +\mu_{h}^{-1}(\nu_h(\sigma_h(x))),
\end{align} 
which belongs to $I$. 

For $(h, g) \in \Lambda_H$ and $(x, y) \in I \times I$,  using \eqref{action1} at the forth equality, we make the following computation  in $\Lambda_E$: 
\begin{align*}
(s(h), s(g))^{-1} (x, y) (s(h), s(g)) &= \big(\lambda_{s(g)^{-1}}(- s(h)), {s(g)^{-1}}\big)(x, y) (s(h), s(g))\nonumber\\
&= \big(\lambda_{{s(g)^{-1}}}(- s(h)) + \lambda_{{s(g)^{-1}}}(x) + \lambda_{{s(g)^{-1}} \circ y} (s(h)), {s(g)^{-1}} \circ y \circ s(g)\big)\nonumber\\
&= \lambda_{{s(g)^{-1}}}\big(-s(h) + x + \lambda_y(s(h)), \sigma_g(y)\big)\nonumber\\
&=  \lambda_{{s(g)^{-1}}}\big(-s(h) + x -y + \mu^{-1}_h(\nu_h(\sigma_h(y))), \sigma_g(y)\big)\nonumber\\
&=  \big(\nu^{-1}_g\big(\mu_h(x - y) + \nu_h(\sigma_h(y))\big),  \sigma_g(y)\big)\nonumber\\
&=  \big(\nu^{-1}_g(\mu_h(x - y)) + \nu^{-1}_g(\nu_h(\sigma_h(y))),  \sigma_g(y)\big).
\end{align*}
Here comes the action $\chi : \Lambda_H \to \Aut(I \times I)$; $(h, g)  \mapsto \chi_{(h, g)}$, where  
\begin{equation}\label{Lambda-action}
\chi_{(h,g)}(x,y) = \big(\nu^{-1}_g(\mu_h(x - y)) + \nu^{-1}_g(\nu_h(\sigma_h(y))),  \sigma_g(y)\big).
\end{equation}
We remark that the action $\chi$ of $\Lambda_H$ on $I \times I$ depends on the good triplet of actions  $(\nu, \mu, \sigma)$, and not on the choice of the extension $E$. Also, if $I$ is chosen so that $\lambda_x(w) = w$ for all $x \in I$ and $w \in E$, then $\chi$ takes an easy form
\begin{equation}\label{ker-action}
\chi_{(h,g)}(x,y) = \big(\nu^{-1}_g(\mu_h(x ),  \sigma_g(y)\big).
\end{equation}
Indeed, if $\lambda_x(w) = w$ for all $x \in I$ and $w \in E$, then, for all $h \in H$ and $y \in I$, we get 
\begin{align*}
	\nu_h (\sigma_h(y))& =\nu_h(s(h)^{-1} \circ y \circ s(h) )= -s(h)+ (y \circ s(h)) = -s(h)+ y + \lambda_y(s(h)) \\
&=  -s(h)+ y +s(h) =\mu_h(y).
\end{align*}



The following result explains how compatibility conditions of a given arbitrary triple of actions $(\nu, \mu, \sigma)$ is connected to the group action of $\Lambda_H$ on $I \times I$ given by \eqref{Lambda-action}.

\begin{thm}
	Let $H$ be a skew brace and $I$ an abelian group. Let
	\[
	\mu: (H, +) \rightarrow \Aut(I) \quad \text{and} \quad \sigma: (H, \circ) \rightarrow \Aut(I)
	\]
	be antihomomorphisms, and let $\nu: (H, \circ) \rightarrow \Aut(I)$ be a homomorphism. Define a map
	\[
	\chi^{(\nu, \mu, \sigma)} : \Lambda_H \longrightarrow \Aut(I \times I)
	\]
	by setting
	\[
	\chi^{(\nu, \mu, \sigma)}_{(h,g)}(x,y) = \left(  \nu^{-1}_g(\mu_h(x - y)) + \nu^{-1}_g(\nu_h(\sigma_h(y))),  \sigma_g(y)\right)
	\]
	for all $h, g \in H$ and $x, y \in I$. Then $\chi^{(\nu, \mu, \sigma)}$ is an anti-homomorphism if and only if $(\nu, \mu, \sigma)$ is a good triplet of actions.
\end{thm}
\begin{proof}
Suppose that $\chi^{(\nu, \mu, \sigma)}$ is an anti-homomorphism. Then it is easy to see that the subgroup $I \times \{0\}$ is invariant under the action of  $\chi^{(\nu, \mu, \sigma)}$. 
Indeed, for any $h, g \in H$ and $x \in I$, we have
\[
\chi^{(\nu, \mu, \sigma)}_{(h,g)}(x,0) = \left( \nu^{-1}_g(\mu_h(x)) , \; 0\right).
\]
Note that $(0,g)(h,0) = (\lambda_g(h), g)$. Using this and the fact that $\chi^{(\nu, \mu, \sigma)}$ is an anti-homomorphism, we obtain
\begin{align*}
	\left( \mu_h(\nu^{-1}_g(x)), \; 0 \right)  &=  \chi^{(\nu, \mu, \sigma)}_{(h,\;0)}
	 \chi^{(\nu, \mu, \sigma)}_{(0,\;g)}(x,\;0)  = 	\chi^{(\nu, \mu, \sigma)}_{(0,\;g)(h,\;0)}(x,\;0)\\
	&=   \chi^{(\nu, \mu, \sigma)}_{(\lambda_g(h),\; g)}(x, \;0) = \nu^{-1}_g(\mu_{\lambda_g(h)}(x)),
	\end{align*}
which is the first condition of the good triplet of actions $(\nu, \mu, \sigma)$.


Next, we verify the second compatibility condition. Since $\chi^{(\nu, \mu, \sigma)}$ is an anti-homomorphism, for all $h, g \in H$ and $y \in I$, we have
\begin{align*}
	\chi^{(\nu, \mu, \sigma)}_{(h,0)} \chi^{(\nu, \mu, \sigma)}_{(g,0)}(0,y) 
	= \chi^{(\nu, \mu, \sigma)}_{(g+h,0)}(0,y).
\end{align*}

Making further computations, we obtain
\begin{align} \label{second}
	\chi^{(\nu, \mu, \sigma)}_{(h,0)} \chi^{(\nu, \mu, \sigma)}_{(g,0)}(0,y) 
	&= \chi^{(\nu, \mu, \sigma)}_{(h,0)} \left( -\mu_g(y) + \nu_g(\sigma_g(y)), \; y \right) \notag \\
	&= \left( \mu_h\!\left( -\mu_g(y) + \nu_g(\sigma_g(y)) - y \right) + \nu_h(\sigma_h(y)), \; y \right)
\end{align}
and
\begin{align}\label{third}
	\chi^{(\nu, \mu, \sigma)}_{(g+h,0)}(0,y) 
	&= \left( -\mu_{g+h}(y) + \nu_{g+h}(\sigma_{g+h}(y)), \; y \right).
\end{align}
Comparing the first coordinates, we deduce
\begin{align*}
	\nu_{g+h}(\sigma_{g+h}(y)) 
	= \mu_h\!\big(\nu_g(\sigma_g(y))\big) - \mu_h(y) + \nu_h(\sigma_h(y)),
\end{align*}
which is the second condition of the  good triplet of actions $(\nu, \mu, \sigma)$.

Conversely, suppose that $(\nu, \mu, \sigma)$ is a good triplet of actions. Then there exists a skew brace extension $E$ of $H$ by $I$ with associated actions $(\nu, \mu, \sigma)$. This induces a group extension $\Lambda_E$ of $\Lambda_H$ by $I \times I$. By \eqref{Lambda-action}, the map $\chi^{(\nu, \mu, \sigma)}$ describes the corresponding action of $\Lambda_H$ on $I \times I$. Thus, $\chi^{(\nu, \mu, \sigma)}$ is an anti-homomorphism.
\end{proof}


\begin{cor}\label{inducedaction}
	Let $H$ be a skew brace and let $I$ be an $H$-module via a good triplet of actions $(\nu, \mu, \sigma)$ on $I$. Then we can define an action of $\Lambda_H$ on $I$ by
	\[
	\phi^{(\nu, \mu)}_{(h_1, h_2)} = \mu^{-1}_{h_1} \nu_{h_2}, \quad \text{for all } h_1, h_2 \in H.
	\]
\end{cor}
\begin{proof}
Let $p_1: I \times I \to I$ be the projection at the first factor. Then $\phi^{(\nu, \mu)} :=  p_1 \chi^{(\nu, \mu, \sigma)} : \Lambda_H  \to I$ is the desired action of  $\Lambda_H$ on $I$, where $\chi^{(\nu, \mu, \sigma)}$ is defined in the preceding result. \end{proof}

\begin{remark}\label{remark-ideal}
Let $G$ be a skew brace and $I$ be its ideal. Then equation \eqref{action1} also holds in more generality, that is, $\lambda_x(g)-g \in I$ for all $x \in I$ and $g \in G$. The proof is essentially the same and uses the fact that $I$ is normal in both the groups $(G, +)$ and $(G, \circ)$.
\end{remark}


One might expect that a non-split brace extension $0 \to I \to E \to H \to 0$ always goes to a non-split extension $0 \to \Lambda_I \to \Lambda_E \to \Lambda_H \to  0$. But this is not the case, as explained in the following example.  
\begin{example}\label{exp0}
Let $A$ be the brace with $(A, +) = \mathbb{Z}_2 \times \mathbb{Z}_4$, whose socle has order $2$, and $(A, \circ)$ defined as follows:
\[
(x_1, y_1) \circ (x_2, y_2)
\;:=\;
\bigl(x_1 + x_2,\; y_1 + y_2 + 2y_1x_2 + 2(x_1 + y_1)y_2 \bigr).
\]
In this case, as classified in \cite{Bach15}, we get
\[
(A, \circ) \;\cong\; (\mathbb{Z}/2)^3.
\]
It is easy to see that 
\[
I = \mathbb{Z}_2 \times \langle \bar{2} \rangle \;\cong\; \mathbb{Z}_2 \times \mathbb{Z}_2
\]
is an ideal of $A$, which is a trivial brace. The GAP id of $A$ is \texttt{SmallSkewbrace(8,8)}. A computation by GAP \cite{GAP}, using the package Yang-Baxter \cite{VK24}, shows that the ideal $I$ does not admit a complement in $A$, but $\Lambda_I$ does admit a complement in $\Lambda_A$. Hence, the extension $0 \to I \to A \to A/I \to 0$ of  skew braces does not split, but the group extension $0 \to \Lambda_I \to \Lambda_A \to \Lambda_A/\Lambda_I \to 0$  does split.
\end{example}
 
So, a natural question, which arises here, is: Are there certain conditions under which  non-split brace extensions  always go to associated  non-split group extensions?  The answer is affirmative, as we see below. Let a skew brace $H$ act on an abelian group $I$ via a good triplet of actions $(\nu, \mu, \sigma)$, and $E$ be the associated brace extension. Then  $E$ is  a  socle extension of $H$ by $I$ if and only if $\mu$ is a trivial action of $(H, +)$ on $I$ and $\nu_h^{-1} = \sigma_h$ for all $h \in H$.  Indeed,
if $I \le \Soc(E)$, then $\mu$ is trivial, and using  \eqref{action1}, that is,
\[
0 = \lambda_x(s(h)) - s(h) = -x + \mu_{h}^{-1}\bigl(\nu_h(\sigma_h(x))\bigr),
\]  
we obtain $\nu_h(\sigma_h(x)) = x$ for all $x \in I$.  
For the converse, if $\mu$ is trivial, then $I \subseteq \Z(E,+)$.  
To show that $I \subseteq \ker \lambda$, it is enough to check that $\lambda_x(s(h)) = s(h)$ for all $h \in H$ and $x \in I$, which easily follows from \eqref{action1}.  We now answer the above question.

\begin{thm}\label{prop-split-ext}
	The skew brace extension \eqref{sb-ext}  with $I \le \Soc(E)$ splits if and only if  the associated group extension \eqref{g-ext} splits.
\end{thm}

\begin{proof}
Note that `only if' part of the assertion follows from the construction.  So assume that the group extension \eqref{g-ext} splits. Then there exists a section (group homomorphism) 
$s : \Lambda_H \to \Lambda_E$ such that  \( (\pi \times \pi) \circ s = \mathrm{Id}_{\Lambda_H} \). For all \( h, g \in H \), we can write
$$
s(h, g) = (s_1(h, g), s_2(h, g)),$$
where \( s_1, s_2 : H \times H \to E \) are the component maps of $s$.
Since \( s \) is a group homomorphism, for all \( (h_1, g_1), (h_2, g_2) \in \Lambda_H \), we have
$$s(h_1 + \lambda_{g_1}(h_2),\; g_1 \circ g_2) = s(h_1, g_1) \cdot s(h_2, g_2).$$
Now writing in the  component form and making further computations we get
$$s(h_1 + \lambda_{g_1}(h_2),\; g_1 \circ g_2) = (s_1(h_1 + \lambda_{g_1}(h_2),\; g_1 \circ g_2),\; s_2(h_1 + \lambda_{g_1}(h_2),\; g_1 \circ g_2)) $$
and 
\begin{align*}
s(h_1, g_1) \cdot s(h_2, g_2)  &= (s_1(h_1, g_1), s_2(h_1, g_1)) \cdot (s_1(h_2, g_2), s_2(h_2, g_2)) \\
		&= \left(s_1(h_1, g_1) + \lambda_{s_2(h_1, g_1)}(s_1(h_2, g_2)),\; s_2(h_1, g_1) \circ s_2(h_2, g_2)\right).
	\end{align*}
	Comparing the first components, we obtain the key relation:
	\begin{equation}\label{eq:main1}
		s_1(h_1 + \lambda_{g_1}(h_2),\; g_1 \circ g_2) = s_1(h_1, g_1) + \lambda_{s_2(h_1, g_1)}(s_1(h_2, g_2)).
	\end{equation}
	
	Now, $s$ being a section, we see that \( \pi(s_1(h, g)) = h \) and \( \pi(s_2(h, g)) = g \), and hence the difference \( s_2(h, g) - s_1(g, h) \in I \). So there exists a  map \( f : H \times H \to I \) such that 
		\[
	s_2(h, g) = s_1(g, h) + f(h, g).
	\]
 In particular, $s_1(0, h)$ and $s_2(h, 0)$ both lie in $I$ for all $h \in H$. Observe that for any $h_1, h_2 \in H$, we have
	\[
	s(h_1, h_2) = s((h_1, 0)(0, h_2)),
	\]
	which, using \eqref{eq:main1},  implies that
	\begin{equation}\label{eq:addsplit}
		s_1(h_1, h_2) = s_1(h_1, 0) + s_1(0, h_2),
	\end{equation}
	since $s_2(h_1, 0) \in I$.

	
Using the fact that $I \le \Soc(E)$ and the identity $s_2(h, g) = s_1(g, h) + f(h, g)$, \eqref{eq:main1} takes the form
	\begin{equation}\label{eq:main2}
		s_1(h_1 + \lambda_{g_1}(h_2),\; g_1 \circ g_2) = s_1(h_1, g_1) + \lambda_{s_1(g_1, h_1)}(s_1(h_2, g_2)).
	\end{equation}
Define a map \( t : H \to E \) by
	\[
	t(h) := s_1(h, 0).
	\]
	That $t$ is a st-section of the projection \( \pi : E \to H \) is obvious. We claim that \( t \) is a skew brace homomorphism. Taking  \( g_1 = g_2 = 0 \) in  \eqref{eq:main2}, we get
	$$s_1(h_1 + h_2, 0) = s_1(h_1, 0) + \lambda_{s_1(0, h_1)}(s_1(h_2, 0))$$
	for all  \( h_1, h_2 \in H \). 	Since \( s_1(0, h_1) \in I \), this, in the form of $t$, gives
	\[
	t(h_1 + h_2) = t(h_1) + t(h_2).
	\]
This shows that \( t \) is a homomorphism of the additive groups. Next, taking  \( h_1 = 0 \) and \( g_2 = 0 \) in \eqref{eq:main2}, we get
$$s_1(\lambda_{g_1}(h_2), g_1) = s_1(0, g_1) + \lambda_{s_1(g_1, 0)}(s_1(h_2, 0))$$
for all $g_1, h_2 \in H$, which,
	using  \eqref{eq:addsplit} on the left hand side and the fact $I \le \Soc(E)$, gives
	\[
	s_1(\lambda_{g_1}(h_2), 0) = \lambda_{s_1(g_1, 0)}(s_1(h_2, 0)).
	\]
	This, using the definition of $t$, gives 
	\[
	t(\lambda_{g_1}(h_2)) = \lambda_{t(g_1)}(t(h_2)).
	\]
Since $t$ is  a homomorphism of the additive groups, we get $t(g_1 \circ h_2) = t(g_1) \circ t(h_2)$ for all $g_1, h_2 \in H$. Hence $t$ is a skew brace homomorphism, and the proof is complete. 		
\end{proof}


We now explore a connection between the second cohomology group of  a skew brace $H$ with coefficients in an abelian group $I$ having $(\nu, \mu, \sigma)$ as a good triplet of actions with the second cohomology group of $\Lambda_H$ with coefficients in  $I \times I$ having  induced action $\chi^{(\nu, \mu, \sigma)}$. 
Let \( E \) be a brace extension of   \( H \) by  \( I \),  and let \( s: H \rightarrow E \) be an st-section. For \( h_1, h_2 \in H \), the associated \( 2 \)-cocycles of $(H, +)$ and $(H, \circ)$ with coefficients in $I$ are defined by:
\[
\beta(h_1, h_2) := -s(h_1 + h_2) + s(h_1) + s(h_2),
\]
\[
\tau(h_1, h_2) := -s(h_1 \circ h_2) +  (s(h_1) \circ s(h_2)) = \nu_{h_1 \circ h_2}(\overline{\tau}(h_1, h_2)),
\]
where
\[
\overline{\tau}(h_1, h_2) := s(h_1 \circ h_2)^{-1} \circ s(h_1) \circ s(h_2).
\]
Define $f: H \times H \rightarrow I$ by $f(h_1, h_2)=-s(\lambda_{h_1}(h_2))+ \lambda_{s(h_1)}(s(h_2))$. We have 
\begin{align}
-s(\lambda_{h_1}(h_2))+ \lambda_{s(h_1)}(s(h_2))=& -s(-h_1 + (h_1 \circ h_2)) - s(h_1) + (s( h_1) \circ s(h_2))\nonumber\\
=& \beta(-h_1, h_1 \circ h_2)-s(h_1 \circ h_2) -s(-h_1) -s(h_1) + (s( h_1) \circ s(h_2))\nonumber\\
=&  \beta(-h_1, h_1 \circ h_2)-s(h_1 \circ h_2) -\beta(h_1, -h_1) + (s( h_1) \circ s(h_2)) \nonumber\\
=&   \beta(-h_1, h_1 \circ h_2) -\mu_{h_1 \circ h_2}(\beta(h_1, -h_1)) -s(h_1 \circ h_2)+ (s( h_1) \circ s(h_2)) \nonumber\\
=&   \beta(-h_1, h_1 \circ h_2) -\mu_{h_1 \circ h_2}(\beta(h_1, -h_1)) + \tau(h_1, h_2 ). \label{f-map}
\end{align}


Now we compute  $2$-cocycles corresponding to the extension $\Lambda_E$ of $\Lambda_H$ by $\Lambda_I$, with respect to the section $S:=s \times s$. For  $h_1, h_2, g_1, g_2 \in H$, using \eqref{f-map}, we have
\begin{align*}
 &S((h_1,g_1) (h_2, g_2))^{-1} S(h_1, g_1) S(h_2, g_2)\\
&= \big((s \times s)((h_1+ \lambda_{g_1}(h_2), g_1 \circ g_2)) \big)^{-1}(s(h_1), s(g_1) ) (s(h_2), s(g_2))\\
&= \big((s(h_1+ \lambda_{g_1}(h_2)), s(g_1 \circ g_2)) \big)^{-1} \big( s(h_1)+ \lambda_{s(g_1)}(s(h_2)), s(g_1) \circ s(g_2) \big)\\
&= \big( -\lambda^{-1}_{s(g_1 \circ g_2)}(s(h_1+ \lambda_{g_1}(h_2)), s(g_1 \circ g_2)^{-1} \big) \big( s(h_1)+ \lambda_{s(g_1)}(s(h_2)), s(g_1) \circ s(g_2) \big)\\
&=  \big( -\lambda^{-1}_{s(g_1 \circ g_2)}(s(h_1+ \lambda_{g_1}(h_2)) + \lambda^{-1}_{ s(g_1 \circ g_2)}\big(s(h_1)+ \lambda_{s(g_1)}(s(h_2))\big),  s(g_1 \circ g_2)^{-1} \circ s(g_1) \circ s(g_2) \big)\\
&=  \big( \lambda^{-1}_{s(g_1 \circ g_2)}\big(-s(h_1+ \lambda_{g_1}(h_2))+ s(h_1)+ \lambda_{s(g_1)}(s(h_2)) \big), \tilde{\tau}(g_1, g_2)\big)\\
&=   \big( \lambda^{-1}_{s(g_1 \circ g_2)}\big(-s(h_1+ \lambda_{g_1}(h_2))+ s(h_1)+ s(\lambda_{g_1}(h_2))-s(\lambda_{g_1}(h_2))+  \lambda_{s(g_1)}(s(h_2))\big), \bar{\tau}(g_1, g_2)\big)\\
&=  \big( \lambda^{-1}_{s(g_1 \circ g_2)}\big(\beta(h_1, \lambda_{g_1}(h_2))+ f(g_1, h_2)\big),  \bar{\tau}(g_1, g_2)\big)\\
&= \big(\nu^{-1}_{g_1 \circ g_2} \big(\beta(h_1, \lambda_{g_1}(h_2))+ f(g_1, h_2)\big), \bar{\tau}(g_1, g_2) \big).
\end{align*}

Let $\Z^2_{Sb}(H,I)$ denote the set of all brace $2$-cocycles with coefficients in $I$, and $\Z^2_{Gp}((H, +),I)$  and  $\Z^2_{Gp}((H, \circ),I)$ denote the set of group $2$-cocycles of the groups $(H, +)$ and $(H, \circ)$ respectively with coefficients in $I$. Then it follows from \cite[Page 10, Eqn. (13)]{RY24} that $(\beta, \tau) \in \Z^2_{Sb}(H,I)$ if and only if  $\beta \in \Z^2_{Gp}((H, +),I)$ and $\overline{\tau} \in \Z^2_{Gp}((H, \circ),I)$ and  the following equation holds:
\begin{eqnarray}\label{mutual1}
	\mu_{\lambda_{h_1}(h_3)}(\tau(h_1, h_2)) - \tau(h_1, h_2 + h_3) +\tau(h_1, h_3) &=& \nu_{h_1}(\beta(h_2, h_3)) + \mu_{h_1 \circ h_3}(\beta(h_1, - h_1)) \nonumber \\ 
	& & - \; \beta(-h_1, h_1 \circ h_3)   - \; \beta(h_1 \circ h_2, \lambda_{h_1}(h_3)).
\end{eqnarray}


We are now ready to compute desired $2$-cocycles.
\begin{thm}
Let a skew brace $H$ act on a trivial brace $I$ via a good triplet of actions $(\nu,  \mu, \sigma)$ and $\beta, \tau : H \times H \rightarrow I$ be two maps. Define a map
$$(\phi_{(\beta, \tau)}, \overline{\tau}) : \Lambda_H \times \Lambda_H \to I \times I$$
by setting
$$(\phi_{(\beta, \tau)}, \overline{\tau})\big((h_1, g_1), (h_2, g_2)\big)
 =  \big(  \nu^{-1}_{g_1 \circ g_2} \big( \beta(h_1, \lambda_{g_1}(h_2))\big) 
+  f(g_1, h_2),\ \nu^{-1}_{g_1 \circ g_2}(\tau(g_1, g_2)) \big),$$
where $f(g_1, h_2) =  \beta(-g_1, g_1 \circ h_2) - \mu_{g_1 \circ h_2}(\beta(g_1, -g_1))  + \tau(g_1, h_2)$.
Then the map $(\phi_{(\beta, \tau)}, \overline{\tau})$ is a group $2$-cocycle of $\Lambda_H$ with coefficients in $I \times I$ if and only if $(\beta, \tau) \in \Z^2_{Sb}(H,I)$. 
\end{thm}

\begin{proof}
If  $(\beta, \tau) \in Z_{Sb}^2(H, I)$, then we know that there exists a skew brace extension $E$ of $H$ by $I$ and an st-section $s$ of $E$ such that the $2$-cocycle corresponding to $s$ is $(\beta, \tau)$. Hence, it follows from the preceding  construction that the map $(\phi_{\beta, \tau}, \overline{\tau})$ is a group $2$-cocycle of $\Lambda_H$ with coefficients in $I \times I$. 

Conversely, let the map $(\phi_{\beta, \tau},   \overline{\tau})$ be a group $2$-cocycle of $\Lambda_H$ with coefficients in $I \times I$. Then for $(h_i, g_i) \in \Lambda_H$, $1 \le i \le 3$, by the definition of $2$-cocycle,  we get 
\begin{align}\label{bigeqn}
0 &= (\phi_{(\beta, \tau)}, \overline{\tau})((h_2, g_2), (h_3, g_3))- (\phi_{(\beta, \tau)}, \overline{\tau})((h_1+\lambda_{g_1}(h_2), g_1 \circ g_2), (h_3, g_3))\nonumber\\
& \quad +  (\phi_{(\beta, \tau)}, \overline{\tau})((h_1, g_1), (h_2+\lambda_{g_2}(h_3), g_2 \circ g_3)) 
 -\chi_{(-\lambda^{-1}_{g_3}(h_3), \, g^{-1}_3)} ((\phi_{(\beta, \tau)}, \overline{\tau})((h_1, g_1), (h_2, g_2))).
\end{align}
Putting $g_1=g_2=g_3=0$, we get $\beta \in  \Z^2_{Gp}((H, +),I)$ and, similarly, putting $h_1=h_2=h_3=0$, we get $\overline{\tau} \in  \Z^2_{Gp}((H, \circ),I)$. It remains to prove the compatibility condition \eqref{mutual1}. Taking $h_1=g_2=g_3=0$ in \eqref{bigeqn}, by the definition  of $(\phi_{(\beta, \tau)}, \overline{\tau})$, we get the following big equation.
\begin{align*}
& (\phi_{(\beta, \tau)}, \overline{\tau})((h_2, 0), (h_3, 0))- (\phi_{(\beta, \tau)}, \overline{\tau})((\lambda_{g_1}(h_2), g_1 ), (h_3, 0))+  (\phi_{(\beta, \tau)}, \overline{\tau})((0, g_1), (h_2+ h_3,   0)) \nonumber\\
& \quad -\chi_{(-h_3, 0)} ((\phi_{(\beta, \tau)}, \overline{\tau})((0, g_1), (h_2, 0))) \\
& =  (\beta(h_2, h_3), 0)-   \Big( \nu^{-1}_{g_1}\big(\beta(\lambda_{g_1}(h_2), \lambda_{g_1}(h_3)) +\beta(-g_1, g_1 \circ h_3)-\mu_{g_1 \circ h_3}(\beta(g_1, -g_1))+ \tau(g_1, h_3)\big), 0\Big)\\
&\quad  + \Big(  \nu^{-1}_{g_1}\big(\beta(-g_1, g_1 \circ(h_2+h_3) )-\mu_{g_1 \circ(h_2+h_3)}(\beta(g_1, -g_1)) + \tau(g_1, h_2+h_3)\big), 0 \Big)\\
& \quad -  \Big(\mu_{h_3}\big(\nu^{-1}_{g_1}\big(\beta(-g_1, g_1 \circ h_2)
-\mu_{g_1 \circ h_2}(\beta(g_1, -g_1))+ \tau(g_1, h_2)\big)\big), 0 \Big) 
\end{align*}

We now re-interpret some parts of the preceding equation to simplify it. Using the first compatibility condition of good triplet of actions, i.e., $\mu_{\lambda_{h_1}(h_3)}(y) = \nu_{h_1} ( \mu_{h_3} ( \nu_{h_1}^{-1}(y) ))
$,  we get
\begin{align*}
\mu_{h_3}\big(\nu^{-1}_{g_1}\big(\beta(-g_1, g_1 \circ h_2) -\mu_{g_1 \circ h_2}(\beta(g_1, -g_1))+ \tau(g_1, h_2)\big) \big) &= \nu^{-1}_{g_1} (\mu_{\lambda_{g_1}(h_3)}(\beta(-g_1, g_1 \circ h_2)\\
&  \quad -\mu_{g_1 \circ h_2}(\beta(g_1, -g_1))+ \tau(g_1, h_2))).
\end{align*}
Also, using the $2$-cocycle property of $\beta$, we get 
\begin{align*}
\beta(-g_1, g_1 \circ(h_2+h_3)) & = \beta(-g_1, g_1 + \lambda_{g_1}(h_2+h_3))\\
& = -\beta(g_1, \lambda_{g_1}(h_2+h_3))+\mu_{\lambda_{g_1}(h_2+h_3)}(\beta(-g_1, g_1))\\
& = \beta(\lambda_{g_1}(h_2), \lambda_{g_1}(h_3))-\beta(g_1 \circ h_2, \lambda_{g_1}(h_3))-\mu_{\lambda_{g_1}(h_3)}(\beta(g_1, \lambda_{g_1}(h_2)))\\
& \quad +\mu_{\lambda_{g_1}(h_2+h_3)}(\beta(-g_1, g_1))
\end{align*}
and 
$$\beta(-g_1, g_1+\lambda_{g_1}(h_2))= -\beta(g_1, \lambda_{g_1}(h_2))+\mu_{\lambda_{g_1}(h_2)}( \beta(-g_1, g_1)).$$

Using these expressions in the big equation and equating to zero component-wise, we finally get 
\begin{align*}
0 &= \beta(h_2, h_3)+ \nu^{-1}_{g_1}\big( -\beta(-g_1, g_1 \circ h_3)+ \mu_{g_1 \circ h_3}(\beta(g_1,-g_1))\\ 
& \quad  -\tau(g_1, h_3)-\mu_{\lambda_{g_1}(h_3)}(\tau(g_1, h_2)) + \tau(g_1, h_2+h_3)\big),
\end{align*}
which is equivalent to the required compatibility condition \eqref{mutual1}, and the proof is complete. 
\end{proof}

We remark that the map $(\beta, \tau) \mapsto (\phi_{(\beta, \tau)}, \overline{\tau})$ from $\Z^2_{Sb}(H,I)$ to 
$\Z^2_{Gp}(\Lambda_H,I \times I)$ is an injective map, and therefore an injective group homomorphism. Indeed, if  $(\phi_{(\beta, \tau)}, \overline{\tau})$ is a zero map, then $\tau(g_1, g_2) = 0$ for all $g_1, g_2 \in H$, and 
$$ \nu^{-1}_{g_1 \circ g_2} \big( \beta(h_1, \lambda_{g_1}(h_2)) = f(g_1, h_2) =  \beta(-g_1, g_1 \circ h_2) -\mu_{g_1 \circ h_2}(\beta(g_1, -g_1))$$
for all $h_1, h_2, g_1, g_2 \in H$. Substituting $g_1 = 0$, we get $\beta(h_1, h_2) = 0$. Hence $(\beta, \tau) = 0$. Since $I \times I$ is an abelian groups, the association is a group homomorphism too. More precisely, one can  easily  verify that
	\[
	\phi_{(\beta_1 + \beta_2, \tau_1 + \tau_2)} = \phi_{(\beta_1, \tau_1)} + \phi_{(\beta_2, \tau_2)}
	\quad \text{and} \quad 
	\overline{\tau_1 + \tau_2} = \overline{\tau_1} + \overline{\tau_2}.
	\]

We now prove that this map induces a group homomorphism at the cohomology level. But we do not know whether this induced group homomorphism is an embedding. However, as we see below, this homomorphism is an embedding in a special case  of actions. 
\begin{thm}\label{Psi-hom}
Let $H$ be a skew brace acting on a trivial brace $I$ via the actions $(\nu, \mu, \sigma)$. Then the  map $\Psi: \Ho^2_{Sb}(H, I) \to \Ho^2_{Gp}(\Lambda_H, I \times I)$, defined by
$$\Psi([(\beta, \tau)]) = [(\phi_{\beta, \tau}, \overline{\tau})]$$
is a group homomorphism.
\end{thm}
\begin{proof}
That 	the map $\Psi$ is well-defined, follows from the fact that the corresponding map
	\[
	\psi: \Ext_{(\nu, \mu, \sigma)}(H, I) \rightarrow \Ext_{\chi^{(\nu, \mu, \sigma)}}(\Lambda_H, I \times I);
	\quad [E] \mapsto [\Lambda_E]
	\]
on the set of equivalence classes of extensions is  well-defined. Indeed, as we know that there exists a bijection between the second cohomology groups and the set of equivalence classes of extensions. So the map $\Psi$ is induced by $\psi$ at the level of cohomology groups. It follows from the preceding discussion that $\Psi$ is a group homomorphism, and the proof is complete.
\end{proof}


For some specific  actions, using Theorem \ref{prop-split-ext} the preceding theorem  gives the following useful result.
\begin{cor}\label{useful-cor}
Let $H$ be a skew brace acting on an abelian group $I$ via a good triplet of actions $(\nu, \mu, \sigma)$ with $\nu_h^{-1} = \sigma_h$ and $\mu_h = \Id$ for all $h \in H$. Then the map
	\[
	\Psi: \Ho^2_{Sb}(H, I) \rightarrow \Ho^2_{Gp}(\Lambda_H, I \times I)
	\]
	of Theorem \ref{Psi-hom} is an embedding.
\end{cor}


To defend our statement about the usefulness of the preceding result, we derive \cite[Theorem 2.2]{LV24}. 


 \begin{prop}
The exponent of the Schur multiplier of a finite skew brace $H$ divides $|H|^2$.
\end{prop}
\begin{proof}
Recall  that  if $G$ is a group and $I_1$ and $I_2$ are $G$-modules, then
$$\Ho^2_{Gp}(G, I_1 \oplus I_2) \cong \Ho^2_{Gp}(G, I_1) \oplus \Ho^2_{Gp}(G, I_2).$$
By the preceding result, by taking $I = \mathbb{C}^*$, the group $\Ho^2_{Sb}(H, \mathbb{C}^*)$, the Schur multiplier of the skew brace $H$, embeds into
$$\Ho^2_{Gp}(\Lambda_H, \mathbb{C}^* \times \mathbb{C}^*) \cong \Ho^2_{Gp}(\Lambda_H, \mathbb{C}^*) \oplus \Ho^2_{Gp}(\Lambda_H, \mathbb{C}^*).$$
It is a known group theory result that the exponent of $\Ho^2_{Gp}(\Lambda_H, \mathbb{C}^*)$ divides $|\Lambda_H|^2$. Hence the exponent of $\Ho^2_{Sb}(H, \mathbb{C}^*)$  divides $|H|^2$, which completes the proof.
\end{proof}

We now present some other applications of Corollary \ref{useful-cor}.  For simplicity of notation, let $M_b(H)$ denote the Schur multiplier of a skew brace $H$, and let $M(G)$ denote the Schur multiplier of a group $G$. We first observe 

\begin{prop}
	Let $H$ be a finite trivial or almost trivial skew brace. Then the exponent of the Schur multiplier of the skew brace $H$ divides the exponent of $H$.
\end{prop}
\begin{proof}
	If $H$ is a trivial or almost trivial skew brace, then $\Lambda_H = H \times H$. 
	It is well known that 
	\[
	\Ho^2_{Gp}(H \times H, \mathbb{C}^*) 
	= \Ho^2_{Gp}(H, \mathbb{C}^*) \oplus \Ho^2_{Gp}(H, \mathbb{C}^*) \oplus \bigl(H/\gamma_2(H) \otimes H/\gamma_2(H)\bigr)
	\]
	(see \cite{GK87}). Moreover, the exponent of each factor divides the exponent of $H$. The assertion now  follows using Corollary \ref{useful-cor}.
\end{proof}

The preceding result also follows from  the following much stronger result  proved in \cite[Proposition 2.4]{LV24}.

\begin{thm}
	Let $G$ be a group. Then $M_b(G) \cong M_b(G^{\mathrm{op}}) \cong M(G) \times (G/\gamma_2(G) \otimes G/\gamma_2(G))$.
\end{thm}

Another application to braces of prime power orders is the following.
\begin{prop}
	Let $(H, +, \circ)$ be a finite skew brace of abelian type having order $p^n$, where $p$ is an odd prime. Then the exponent of $M_b(H)$ divides $|H|$.
\end{prop}

\begin{proof}
	Let $(H, +, \circ)$ be a brace of order $p^n$, with $p$ an odd prime. It follows from  \cite[Section 2]{Evens72} that
	\[
	M(\Lambda_H) \cong M(H, \circ) \oplus H_1((H, \circ), (H, +)) \oplus M(H, +)_{(H, \circ)},
	\]
	where $M(H, +)_{(H, \circ)}$ is the maximal $(H, \circ)$-trivial quotient of $M(H, +)$, defined as
	\[
	M(H, +)_{(H, \circ)} = M(H, +) / \langle h \cdot x - x \mid h \in H, x \in M(H, +) \rangle^+,
	\]
	with `$\cdot$' denoting the action of $(H, \circ)$ on $M(H, +)$. 
	
	It is straightforward to see that the exponent of each factor divides $|H|$. Hence, using Corollary \ref{useful-cor}, the exponent of $M_b(H)$ divides $|H|$, and the proof is complete.
\end{proof}


We can define a homomorphism from $\Ho^2_{Sb}(H, I)$ to  $\Ho^2_{Gp}((H, \circ), I)$; $[(\beta, \tau)]\mapsto [\overline{\tau}]$. It is not difficult to show that $\bar{\tau} = \nu^{-1}_{h_1 \circ h_2}\tau$ is a group $2$-cocycle of $(H, \circ)$ with coefficients in $I$ and action $\sigma$. If $[(\beta, \tau)]\mapsto [\overline{\tau}]$ is well defined, then, being the projection map, it will a group homomorphism. So we only need to show that  the map is well defined, that is, if $(\beta_1, \tau_1)$ and $(\beta_2, \tau_2)$ are cohomologous, then so are the group $2$-cocycles $\bar{\tau}_1$ and $\bar{\tau}_2$. Assume that $(\beta_1, \tau_1)$ and $(\beta_2, \tau_2)$ are cohomologous. Then there exists a map $\theta: H \to I$ such that (see  \cite[proof of Theorem 3.6]{RY24})
$$(\tau_1 - \tau_2)(h_1, h_2) = \nu_{h_1}(\theta(h_2)) - \theta(h_1 \circ h_2) + \nu_{h_1 \circ h_2}(\sigma_{h_2}(\nu^{-1}_{h_1}(\theta(h_1)))).$$
For $h_1, h_2 \in H$, we now compute
\begin{align*}
(\bar{\tau}_1 - \bar{\tau}_2)(h_1, h_2) &= \nu^{-1}_{h_1 \circ h_2} \left( (\tau_1 - \tau_2)(h_1, h_2) \right)\\
	&= \nu^{-1}_{h_1 \circ h_2} \big( \nu_{h_1}(\theta(h_2)) - \theta(h_1 \circ h_2) + \nu_{h_1 \circ h_2}(\sigma_{h_2}(\nu^{-1}_{h_1}(\theta(h_1)))) \big)\\
	&= \theta_1(h_2) - \theta_1(h_1 \circ h_2) + \sigma_{h_2}(\theta_1(h_1)),
\end{align*}
where $\theta_1(h) = \nu^{-1}_h(\theta(h))$. Hence, $\bar{\tau}_1$ and $\bar{\tau}_2$ are cohomologous via the map $\theta_1$.




We now make some remarks on some specific group actions. The right action of $\Lambda_H$ on $I \times I$ is given by
	\[
	\chi_{(h,g)}(x,y) = \big(\nu_g^{-1} \mu_h(x - y) + \nu_g^{-1} \nu_h
	\sigma_h(y), \sigma_h(y)\big)
	\]
 restricts  to $\Diag(\Lambda_H)$ to give an action on $J := \Lambda_I / \Diag(\Lambda_I)$
such that
	\[
	\chi_{(g,g)}(\overline{(x, y)}) = \overline{(\nu_g^{-1} \mu_g(x - y), 0)}.
	\]
	This action of $\Diag(\Lambda_H)$ is equivalent to the action of $\Diag(\Lambda_H)$ on $I \times \{0\}$; that is, $I \times \{0\}$ is isomorphic to $J$ as a $\Diag(\Lambda_H)$-module under the map
	\[
	(x, 0) \mapsto \overline{(x, 0)}.
	\]
We know that $(H, \circ)$ acts on $I$ via $\sigma$. Note that $\Diag(H) \cong (H, \circ)$; $(h, h) \mapsto h$ and $I \cong I \times \{0\} \cong J$. 

It then follows from  \cite[Proposition 2.2] {PSY18} that the second cohomology groups 
$$\Ho^2_{Gp}\big(\Diag(\Lambda_H), J \big)  \cong \Ho^2_{Gp}\big((H, \circ), I\big)$$
as the given  isomorphisms are action compatible. 





We conclude this section with some investigations on extensions and second cohomology group of the opposite skew brace of a given skew brace $H$ acting non-trivially on an abelian group $I$.
 For a skew brace $E := (E, +, \circ)$, we define the opposite skew brace  $E^{\op} := (E, +^{\op}, \circ)$. We first make a general remark for non-abelian group $I$. If $E$ is an extension of $H$ by $I$, then $E^{\op}$ is an extension of $H^{\op}$ by $I^{\op}$, as $(E/I)^{\op} \cong E^{\op}/I^{\op}$ under the map $x + I \mapsto x +^{\op} I^{\op}$ for all $x \in E$. Let $\Ext(X, I)$ denote the set of equivalence classes of all skew brace extensions of a skew brace $X$ by $I$. Let $\varphi: \Ext(H, I) \rightarrow \Ext(H^{\op}, I^{\op})$ be a map given by   $[E] \mapsto [E^{\op}]$. It follows from  a straightforward computation that $\varphi$ is a well-defined bijective map. 

\begin{prop}\label{bijection1}
The map $\varphi$ is a well-defined bijective map.
\end{prop} 



Now let $I$ be a trivial brace such that a skew brace $H$ acts on it via  $(\nu, \mu, \sigma)$, and $E$ be an extension of $H$ by $I$
\[
0 \longrightarrow I \xrightarrow{i} E \xrightarrow{\pi} H \longrightarrow 0.
\]
Let 
\[
0 \longrightarrow I \xrightarrow{i} E^{\op} \xrightarrow{\pi} H^{\op} \longrightarrow 0,
\]
be the corresponding opposite extension. We determine the action associated with $E^{\op}$.
Let $s$ be an st-section of $E$. Then $s$ is also a section of $E^{\op}$. For $h \in H^{\op}$ and $y \in I$, we have
$$\mu^{\op}_h(y) = -s(h) +^{\op} y +^{\op} s(h) = s(h) + y - s(h) = \mu_{-h}(y).$$
Similarly, we compute
\begin{align*}
	\nu^{\op}_h(y) 
	&= -s(h) +^{\op} (s(h) \circ y) = (s(h) \circ y) - s(h) \\
	&= s(h) - s(h) + (s(h) \circ y) - s(h) = \mu_{-h}(\nu_h(y)).
\end{align*}
Moreover, $\sigma^{\op}_h = \sigma_h$. Hence, this shows that the map $\varphi$ defined in Proposition~\ref{bijection1} restricts to a bijection
\[
\varphi : \Ext_{(\nu, \mu, \sigma)}(H, I) \longrightarrow \Ext_{(\mu^{\op}, \nu^{\op}, \sigma^{\op})}(H^{\op}, I).
\]

%\begin{remark}
%The action $\nu^{\op}$ of $H^{\op}$ on $I$ is the same as the action of diagonal subgroup of  $\Lambda_H$ restricted to $I \times \{0\}$.
%\end{remark}

Next we explore the connection between second cohomology groups of $H$ and $H^{\op}$ with coefficients in a trivial brace $I$. For $h_1, h_2 \in H$, we have
\begin{align*}
\beta^{\op}(h_1, h_2) & = -s(h_1+^{\op}h_2)+^{\op} s(h_1) +^{\op}s(h_2)\\
& =s(h_2)+s(h_1) -s(h_2+h_1)\\
& =s(h_2)+s(h_1) -s(h_2+h_1)+s(h_2)+s(h_1)-s(h_1)-s(h_2)\\
&= \mu_{-(h_2+h_1)}(\beta(h_2, h_1))
\end{align*}
and 
\begin{align*}
\tau^{\op}(h_1, h_2) & =-s(h_1 \circ h_2)+^{\op} (s(h_1) \circ s(h_2))\\
&=(s(h_1) \circ s(h_2)) - s(h_1 \circ h_2)\\
&= (s(h_1) \circ s(h_2)) -s(h_1 \circ h_2) + (s(h_1) \circ s(h_2)) - (s(h_1) \circ s(h_2))\\
&= \mu_{-(h_1 \circ h_2)}(\tau(h_1, h_2)), \mbox{ (since }s(h_1) \circ s(h_2) - s(h_1 \circ h_2) \in I).
\end{align*}
Putting all together we get the following result, which generalizes \cite[Theorem 1.37]{LV24}.
\begin{thm}
Let $H$ be a skew  brace acting on a  trivial brace $I$ via a good triplet of actions $(\nu, \mu, \sigma)$. Then the following hold true:
\begin{itemize}
\item[(i)] The pair $(\beta, \tau) \in \Z^2_{Sb}(H,I)$ if and only if  $(\beta^{\op}, \tau^{\op}) \in \Z^2_{Sb}(H^{\op},I)$.
\item[(ii)] The map $\tilde{\varphi}: \Ho^2_{Sb}(H, I) \rightarrow \Ho^2_{Sb}(H^{\op}, I)$ given by $[(\beta, \tau)] \mapsto [(\beta^{\op}, \tau^{\op})]$
is an isomorphism.
\end{itemize}
\end{thm}


%%%%%%%%%%%%%%%%%%%%%%%%%%%%%%%%%%%%%%%%%%%%%

\section{Complements and supplements}

In this section we first get several fundamental results on complements and supplements of ideals in skew braces. We then obtain results on Frattini sub-skew brace of a skew brace and minimal extensions of skew braces.  Group theoretic analogs of the results on minimal extensions are proved in \cite{Hill72}.  An ideal $I$ of $E$ is said to be  a \emph{Hall ideal} if the orders of $I$ and $E/I$ are coprime. 
\begin{prop}
Let $0 \longrightarrow I \xrightarrow{i} E \xrightarrow{\pi} H \longrightarrow 0$ be a brace extension of a skew brace $H$ by a skew brace $I$. Then the brace extension $E$  splits if and only if there exists a sub-skew brace $G$ of $E$ such that $G+I = E$ and $G \cap I = \{0\}$, that is $I$ admits a complement $G$ in $E$. 
\end{prop}
\begin{proof}
First assume that the brace extension $E$ splits. Then there exists a skew brace homomorphism $s : H \to E$ such that $\pi s = \Id_H$. Then obviously $s(H)$ is the required $G$.
Conversely, note  that the map $\pi|_G: G \rightarrow H$ is an isomorphism of skew braces. Define $s: H \rightarrow E$ by $s(h) = (\pi|_G)^{-1}(h)$. Then $s$ is a skew brace homomorphism and  a section of $\pi$. Hence $E$ is a split extension. 
\end{proof}

In view of this result, while dealing with natural extension $0 \to I \to E \to G \to 0$, where $I$ is an ideal of the skew brace $E$ and $G = E/I$, we usually say that {\it $E$ splits over $I$} whenever $I$ admits a complement in $E$. 

A finite skew brace $E$ is said to be \emph{soluble} if it  admits an abelian series of ideals, that is, a series of ideals of $E$
	\[
	E = I_n \supseteq I_{n-1} \supseteq \cdots \supseteq I_0 = \{0\},
	\]
such that $I_{i}/I_{i-1}$ is a trivial skew brace of abelian type for every $1 \leq i \leq n$.
The  skew brace $E$ is said to be \emph{supersoluble} if it is soluble with 
each quotient  $I_{i+1}/I_i$  trivial skew brace of prime order.



 A sub-skew brace $K$ of a skew brace $E$ is said to be a Hall $\pi$-sub-skew brace if $K$ is a $\pi$-number and the index of $K$ in $E$ is a $\pi'$-number. If $\pi = \{p\}$, then Hall $\pi$-sub-skew brace $K$ of $E$ is said to be Sylow $p$-sub-skew brace of $E$. We first state the following result  \cite[Theorem~1.6]{CDMFT25}:

\begin{thm}\label{CDMFT}
A finite supersoluble skew brace contains Hall $\pi$-sub-skew braces for every set $\pi$ of primes.
\end{thm}

Our next result is an analog of the Schur-Zassenhaus theorem in several specific cases. We could not prove it in the full generality.

\begin{thm}\label{main7}
	Let $E$ be an extension of $H$ by $I$, where $H$ and $I$ are skew braces of coprime orders.  
	Let $|H| = m$ and $|I| = n$. Then the following statements hold:
	
	\begin{enumerate}
		\item There exist complements $H_1$ of $(I,+)$ in $(E,+)$ and $H_2$ of $(I,\circ)$ in $(E,\circ)$ such that 
		$\lambda_g(h) \in H_1$ for all $g \in H_2$ and $h \in H_1$. Moreover
		\[
		\lambda_x(h) \equiv h \pmod{I}
		\]
	for all  $x \in I$ and $h \in H_1$. 
	
	\item  For any $z \in E$ and $h \in H_1$, there exists an element $g \in H_2$ such that 
		\[
		\lambda_{z}(h) \equiv \lambda_{g}(h) \pmod{I}.
		\]
		
	\item For all $x \in I$ and $h \in H_1$,  $\lambda_x(h) = h$  in any of the following two cases:
	\begin{itemize}
		\item[(i)] $\lambda_x(h)$ commutes with $h$ in $(E, +)$ for all $x \in I$;
		\item[(ii)] $I \le \ker(\lambda)$.
	\end{itemize}
	In these cases, $H_1$ becomes a left ideal of $E$, and therefore a complement of $I$ in $E$. Specifically, if $E$ is of abelian type, then $H_1$ is an ideal of $E$, hence a complement of $I$.
	
	\item Define the map $\phi : H_1 \to I$ by 
	\[
	h = g_h \circ \phi(h)
	\]
	for some $g_h \in H_2$. Then $H_1$ is a sub-skew brace of $E$ if and only if $\lambda_{\phi(h)}(h_1) = h_1$ for all $h, h_1 \in H_1$. Thus,  if $\lambda_{\phi(h)}(h_1) = h_1$ for all $h, h_1 \in H_1$, then $I$ admits a complement in $E$.
	
	\item Let $h \in H_1$ and $g_h \in H_2$ be such that $h = g_h \circ \phi(h)$ for some $\phi(h) \in I$. Then the order of $h$ divides the order of $g_h$ in $(E, +)$ and  the order of $g_h$ divides the order $h$ are  in $(E, \circ)$.

\item If $E$ is supersoluble, then $I$ admits a complement in $E$.

\item If $H_1$ is the only complement of $I$ in $(E, +)$, then $H_1$ is a left ideal of $E$, and therefore a complement of $I$. In particular, $I$ admits a complement if $(E, +)$ is nilpotent.

\end{enumerate}
\end{thm}

\begin{proof}
(1) We know that $\Lambda_I$ is a Hall ideal of $\Lambda_E$. Hence $\Lambda_E$ admits a complement of $\Lambda_I$, say $M$, of order $m^2$.  Let $\pi|_M : M \to (E, \circ)$ be the restriction of the projection map $\pi  : \Lambda_E \to (E, \circ)$ to $M$.  Since $\pi$ is a homomorphism, we have  $\ker(\pi|_M) = M \cap (E,+)$, where $(E, +)$ is viewed as a subgroup of $\Lambda_E$. We, by the first isomorphism theorem of groups, now get
	\[
	|M| = |M \cap (E,+)| \, |\pi|_M(M)|.
	\]
Note that neither $|M \cap (E,+)|$ nor $|\pi|_M(M)|$ can be larger than $m$, which, by the preceding equation,  proves that both of these are equal to $m$. Without loss of any generality, we can assume that $H_2 :=  \pi|_M(M) \cap (E, \circ)$, viewing $(E, \circ)$ as a subgroup of $\Lambda_E$, is a subgroup of  $M$. Indeed, since the order of $H_2$ is $m$, which is coprime to $|\Lambda_I|$, there exists a complement $M'$ of $\Lambda_I$ in $\Lambda_E$ containing $H_2$. So taking $M = M'$, we see that $\pi|_M(M) = H_2$.

Define  $H_1 := M \cap (E,+)$, again viewing $(E, +)$ as a subgroup of $\Lambda_E$.
Since $H_1$ is a normal subgroup of $M$,  we have a natural  exact  sequence $0 \to H_1 \to M \to H_2 \to 0$ of groups, which, for the st-section $b \mapsto (0, b)$ of the projection map $\pi|_M$,  induces  an action of $H_2$ on $H_1$. More precisely, an element $g \in H_2$ acts on  $h \in H_1$ by the conjugation rule 
$$(0, g)^{-1} (h, 0) (0, g) = (\lambda_{g^{-1}}(h), 0).$$
Since $|H_1| = |H_2| = m$, the subgroups $H_1$ and $H_2$ are complements of $I$ in their respective groups. It now  follows that $\lambda_g(h) \in H_1$ for $h_2 \in H_2$ and $h \in H_1$.
Final assertion follows from Remark \ref{remark-ideal}.

(2)  We know, by (1), that for every $z \in E$, there exist $g \in H_2$ and $y \in I$ such that $z=g \circ y$; hence for any $h \in H_1$, we have
 $$\lambda_{z}(h)=	\lambda_{g \circ y}(h ) = \lambda_g(\lambda_y(h) = \lambda_g(h+x) = \lambda_g(h) + \lambda_g(x),$$
 where $x =  \lambda_y(h) \in I$, and therefore $\lambda_g(x) \in I$. The assertion now holds.

 
 (3)  We know by Remark \ref{remark-ideal} that  $\lambda_x(h) - h \in I$ for all $x \in I$ and $h \in H_1$. Since $\lambda_x$ is an automorphism of $(E, +)$,  the orders of $\lambda_x(h)$ and $h$ are divisors of $m$.  Now, as $\lambda_x(h)$ and $h$ commute, it follows that the order of $\lambda_x(h) - h$ in $(E, +)$ also divides $m$. This is possible only when $\lambda_x(h) = h$. The second part  follows trivially. The assertion now holds true by using the fact from (1) that $\lambda_g(h) \in H_1$ for all $g \in H_2$ and $h \in H_1$.
 
 
(4)  Let $H_1$ be a sub-skew brace of $E$. Then $\lambda_{h}(h_1) \in H_1$ for all $h, h_1 \in H_1$.  
Since $h = g_h \circ \phi(h)$ for some $g_h \in H_2$ and $\phi(h) \in I$, we have
\[
\lambda_{h}(h_1) = \lambda_{g_h \circ \phi(h)}(h_1) = \lambda_{g_h}(\lambda_{\phi(h)}(h_1)).
\]
As $\phi(h) \in I$, we can write
\[
\lambda_{\phi(h)}(h_1) = h_1 + x, \quad \text{for some } x \in I.
\]
Thus
\begin{align*}
	\lambda_{h}(h_1) &= \lambda_{g_h}(h_1 + x) \\
	&= \lambda_{g_h}(h_1) + y, \quad \text{where } y =  \lambda_{g_h}(x) \in I.
\end{align*}
We know that  $	\lambda_{h}(h_1),  \lambda_{g_h}(h_1) \in H_1$.    Therefore,
\[
\lambda_{h}(h_1) - \lambda_{g_h}(h_1) \in H_1 \cap I =\{0\},
\]
which gives $y=0$. Since  $\lambda_{g_h}$ is an automorphism of $(E, +)$, $y=0$ implies $x=0$.  Hence $\lambda_{\phi(h)}(h_1) = h_1$.

Conversely, if $\lambda_{\phi(h)}(h_1) = h_1$ for all $h, h_1 \in H_1$, then $\lambda_{h}(h_1) \in H_1$, 
and hence $H_1$ is a sub-skew brace of $E$.

(5)  It is given that $h = g_h \circ \phi(h)$. Then we can write
\[
h = g_h + \lambda_g(\phi(h)).
\]
Let $r$ be the order of $h$ in $(H_1, +)$. Then
\[
2h = g_h + \lambda_{g_h}(\phi(h)) + {g_h} + \lambda_{g_h}(\phi(h)) = 2g_h + f_2(g_h,h)
\]
for some $f_2(g_h,h) \in I$. Proceeding inductively, we obtain
\[
rh = 0 = rg_h + f_r(g_h,h)
\]
for some $f_r(g_h,h) \in I$. Thus $r g_h = -  f_r(g_h,h) \in I$. Since $nh \in H_1$ for all integers $n$, it follows that $ng_h$ can not lie in $I$ for any $n < r$. This shows that the order of $h$ divides the order of $g_h$ in $(E, +)$. 

Now consider $g_h = h \circ \phi(h)^{-1}$. Let $s$ be the order of $g_h$ in $(E, \circ)$. Then 
$$g_h^s =0 =  h^s \circ f_s(h, g_h),$$
where $ f_s(h, g_h) \in I$. Thus $r^s =  f_s(h, g_h)^{-1} \in I$. Since $g_h^n \in H_2$ for all integers $n$, it follows that $h^n$ can not lie in $I$ for any $n < s$. This shows that the order of $g_h$ divides the order of $h$ in $(E, \circ)$. 

(6) This follows from  Theorem~1.6 of \cite{CDMFT25} (see Theorem \ref{CDMFT} above). 

(7)  Let $H_1$ be the unique complement of $(I, +)$ in $(E, +)$.  Now, for any $g \in E$, consider $\lambda_g(H_1)$. Since $H_1 \cap I = \{0\}$, it follows that $\lambda_g(H_1) \cap I = \{0\}$  for all $g \in E$.  This implies that $\lambda_g(H_1)$ is also a complement of $I$ in $(E, +)$ for all $g \in E$. By the uniqueness of the complement, we have $\lambda_g(H_1) = H_1$ for all $g \in E$, which shows that $H_1$ is a left ideal of $E$. Hence $H_1$ is a sub-skew brace of $E$, and the proof is  complete.
\end{proof}



\begin{cor}
Let $H$ be a finite supersoluble skew brace and let $I$ be a finite  abelian group such that 
$|H|$ and $|I|$ are coprime. Then $\Ho^2_{Sb}(H, I)$ is trivial. 
In particular, if $H$ has squarefree order and $\gcd(|H|,|I|)=1$, then 
$\Ho^2_{Sb}(H, I)$ is trivial.
\end{cor}
\begin{proof}
	Since every abelian group, viewed as a trivial brace, is supersoluble, it follows that any extension of $H$ by $I$ is supersoluble. By Theorem~\ref{main7}(6), such an extension must be a semidirect product of $H$ by $I$. Hence, $\Ho^2_{Sb}(H, I)$ is trivial.
\end{proof}

\begin{cor}
	Let $H$ be a  finite skew brace and let $I$ be a finite abelian group such that 
$|H|$ and $|I|$ are coprime. Then  $\Ho^2_{Sb}(H, I)$,  with associated action $(\mu, \nu, \sigma)$,  is trivial if any of the following holds:
	\begin{enumerate}
		\item  $\mu$ is trivial;
		\item $\mu_h(x) = \nu_h(\sigma_h(x))$ for all $h \in H$ and $x \in I$.
	\end{enumerate}
\end{cor}

\begin{proof}
	If $\mu$ is trivial, then every element of $\Ho^2_{Sb}(H, I)$ corresponds to an extension $E$ of $H$ by $I$ such that $I \leq \Z(E, +)$. It is well known that a Hall normal subgroup contained in the center admits a unique complement, which by Theorem \ref{main7}(7) is a left ideal of $E$.  This proves that $I$ admits a complement in $E$, and therefore  $\Ho^2_{Sb}(H, I)$ is trivial.
		
Let  $E$ be an extension of $H$ by $I$. Assume that  $\mu_h(x) = \nu_h(\sigma_h(x))$ for all $h \in H$ and $x \in I$. Using this equation in \eqref{action1}, we immediately get $I \le \Ker(\lambda)$ in $E$. The assertion now follows from  Theorem~\ref{main7}(3).
\end{proof}


\begin{prop}
Let $I$ be  a trivial sub-skew brace  of abelian type of a skew brace $E$, and $G_1$ and $G_2$ be two different complements  of $I$ in $E$ such that $I \le \Ker(\lambda)$, then there exists a brace automorphism of $E$ fixing $I$ and mapping $G_1$ to $G_2$. 
\end{prop}
\begin{proof}
By the given hypothesis, it follows that  for any element $a \in E$, there exists  unique elements $u \in I$ and $g_1 \in G_1$ depending on $a$  such that $a = u + g_1 = u \circ g_1$. Now, for  $g_1 \in G_1$, there exist unique elements $x \in I$ and $g_2 \in G_2$ depending on $g_1$  such that 
$$g_1 = x + g_2 = x \circ g_2.$$ 
By the uniqueness of factorization, it follows that $g_2 = -x + g_1 = x^{-1} \circ g_1$.  Define a map $f: E \to E$ by setting
$$f(a) := u + g_2,$$
where $g_1 = x + g_2$. Then obviously $f$ is well defined, $f(u) =u$ for all $u \in I$ and $f(g_1) = f(0 + g_1) = g_2$. Similarly, we can define $f': E \to E$ by exchanging the roles of $G_1$ and $G_2$, and can prove that $f f' = f' f = \Id_E$. It only remains to prove that $f$ and $f'$ are brace homomorphisms. We only prove it for $f$. Let $a, a' \in E$ with unique factorizations  $a = u + g_1$, $a' = u' + g'_1$, where $u, u' \in I$ and $g_1, g'_1 \in G_1$. Then $a + a' = u + g_1 + u' - g_1 + g_1 + g'_1$. Also $g_1$ and $g'_1$ admit unique factorizations $g_1 = x + g_2$ and $g'_1 = x' + g'_2$, where $x, x' \in I$ and $g_2, g'_2 \in G_2$.  Thus 
\begin{eqnarray*}
f(a + a') &=& u + g_1 + u' - g_1 + g_2 + g'_2\\
&=& u + x + (g_2 + u' - g_2) - x   + g_2 + g'_2\\
&=&  u + g_2 + u' + g'_2 \mbox{ (since $I$ is a trivial brace)}\\
&=& f(a) + f(a').
\end{eqnarray*}
Going verbatim after exchanging `$+$' by `$\circ$', we get $f(a \circ a') = f(a) \circ f(a')$. The proof is complete.
\end{proof}

Next, we state a characterization of the minimal ideal of a finite soluble skew brace, established in \cite[Theorem~B]{BEJP24}, which we'll use in the next result.

\begin{thm}\label{minimumideal}
The minimal ideal of a finite non-zero soluble skew brace is a non-zero trivial skew brace of abelian type. Moreover, it is an elementary abelian \( p \)-group for some prime integer $p$.
\end{thm}

We now present a conditional result on the existence of Sylow $p$-sub-skew brace for every prime divisor $p$ the order of a skew brace $E$.

\begin{thm}
	Let \( E \) be a finite soluble skew brace  such that $\lambda_x(y)$ commutes with  $y$ for all $x, y \in E$.  Then \( E \)  admits a Hall $\pi$-sub-skew brace,  for every subset $\pi$ of the set of all prime divisors of $|E|$.
\end{thm}

\begin{proof}
Let the  order  of $E$ be  \(p_1^{k_1} \cdots p_r^{k_r}\), where $p_i$'s are distinct primes and $k_i$'s and $r$ are natural numbers.  Note that every sub-skew brace and every quotient brace of $E$ satisfy the hypothesis of the theorem.  We proceed by the method of contradiction. So assume that $E$ is  a skew brace of the  smallest order satisfying the given hypothesis, but  not admitting a Hall $\pi'$-sub-skew brace for at least one non-empty subset $\pi'$ of the set $\{p_1, \ldots, p_r\}$.  By Theorem \ref{minimumideal}, $E$ admits a minimal ideal $N$ of prime power order $p^k$ for a prime $p$ and a positive integer $k \ge 1$. So $p = p_i$ for some $i$. Consider the quotient skew brace $E/N$.  Since
 $|E/N| < E|$ and the hypothesis of the theorem holds true for $E/N$,  the assertion holds true for $E/N$.  

Let $\bar{K_{\pi}}$ denote some  Hall $\pi$-sub-skew brace  of $E/N$ for each $\pi$. Let $\{p_i\} \subset \pi''$. If $\pi'' = \{p_i\}$ and  $N$ itself is a Hall $\pi''$-sub-skew brace of $E$, then   we get an exact sequence $0 \to N \to E \to E/N \to 0$. Since $|N|$ is coprime to $|E/N|$, by Theorem \ref{main7}(3), $N$ admits a complement $L$ in $E$. Since $L$ is a proper sub-skew brace  $E$, it admits Hall $\pi$-sub-skew brace, which is also  Hall $\pi$-sub-skew brace of $E$, for all  $\pi$. Now assume that $N$  is not a Hall $\pi''$-sub-skew brace of $E$.  Let $K_{\pi}$ be the pre-image of \( \bar{K_{\pi}} \),  for each $\pi$, in \( E \). Then, obviously, $K_{\pi''}$ is a Hall $\pi''$-sub-skew brace of $E$. For $\pi \ne \pi''$, we get the exact sequence $0 \to N \to K_{\pi} \to \bar{K_{\pi}} \to 0$. Since the order of $N$ is coprime to the order of  $\bar{K_{\pi}}$, by Theorem \ref{main7}(3), $N$ admits a complement $L_{\pi}$ in $K_{\pi}$. It is now clear that $L_{\pi}$ is a Hall $\pi$-sub-skew brace of $E$ for all $\pi \ne \pi''$. This contradicts the minimality of $E$, and the proof is complete.
\end{proof}

As a special case of the preceding theorem, we immediately get
\begin{cor}
	Let \( E \) be a finite soluble skew brace  such that $\lambda_x(y)$ commutes with  $y$ for all $x, y \in E$.  Then \( E \)  admits a Sylow $p$-sub-skew brace for every prime divisor $p$ of $|E|$.
\end{cor}



Next, we examine the structure of supersoluble skew braces of square-free order. As mentioned above, the semidirect product of a skew brace $H$ by a skew brace $I$ will be denoted by $I \rtimes_{Sb} H$. 
First, we recall a basic result from  \cite[Corollary 3.9]{BEFPT24}. 
\begin{thm}
	Every finite skew brace of square-free order is supersoluble.
\end{thm}

\begin{prop} \label{solu1}
	Let $E$ be a finite skew brace of square-free order $n$, and let $n = p_1 p_2 \cdots p_k$ be its prime factorization. 
	Then 
	\[
	E \cong (\cdots((\C_{p_1} \rtimes_{Sb} \C_{p_2}) \rtimes_{Sb} \C_{p_3}) \rtimes_{Sb} \cdots) \rtimes_{Sb} \C_{p_k}.
	\]
\end{prop}

\begin{proof}
	It follows from the preceding theorem that $E$ is a supersoluble skew brace.  Let
	\[
	E = I_k \supseteq I_{k-1} \supseteq \cdots \supseteq I_0 = \{0\},
	\]
	be a  series of $E$ such that each $I_{i}/I_{i-1}$  is a trivial skew brace of prime order. Without loss of generality, assume  that $I_{i}/I_{i-1} \cong \C_{p_i}$ for all $1 \leq i \le k$.  Then $I_1 \cong \C_{p_1}$ and $I_2/I_1 \cong \C_{p_2}$. Since $p_1$ and $p_2$ are coprime, by Theorem~\ref{main7}(6), $I_2$ is a split extension of $\C_{p_2}$ by $\C_{p_1}$. Hence,
	\[
	I_2 \cong \C_{p_1} \rtimes_{Sb} \C_{p_2}.
	\]
	Proceeding inductively, we obtain
	\[
	E \cong (\cdots((\C_{p_1} \rtimes_{Sb} \C_{p_2}) \rtimes_{Sb} \C_{p_3}) \rtimes_{Sb} \cdots) \rtimes_{Sb} \C_{p_k}.
	\]
This completes the proof.
\end{proof}

Next, we recall a description of split extensions of skew braces from \cite[Theorem 3.1, Theorem 3.3]{NR24}.

\begin{thm}\label{splitsbext}
	Let $H$ and $I$ be skew braces, and let $E$ be a split extension of $H$ by $I$.  Then there exist  anti-homomorphisms
	\[
	\mu : (H, +) \to \mathrm{Aut}(I, +), \quad 
	\sigma : (H, \circ) \to \mathrm{Aut}(I, \circ),
	\]
 and a homomorphism $\nu : (H, \circ) \to \mathrm{Aut}(I, +)$
such that $\mu$, $\nu$, and $\sigma$ satisfy a certain compatibility condition.  Then $(E, +, \circ)  \cong I \rtimes_{(\nu, \mu, \sigma)} H$, where  the group operations are given by
	\begin{align}
		(h_1, y_1) + (h_2, y_2) &= (h_1 + h_2,\, \mu_{h_1}(y_2) + y_1), \label{sb+} \\
		(h_1, y_1) \circ (h_2, y_2) &= \big(h_1 \circ h_2,\, 
		\nu_{h_1 \circ h_2}\big(\sigma_{h_2}(\nu_{h_1}^{-1}(y_1)) \circ \nu_{h_2}^{-1}(y_2)\big)\big). \label{sbcirc}
	\end{align}
\end{thm}


We now provide an application of Proposition \ref{solu1}.
\begin{cor}
	Let $E$ be a finite skew brace of square-free order $n$ with prime factorization  $n = p_1 p_2 \cdots p_k$ and $p_i \nmid (p_j - 1)$ for all $1 \leq i, j \leq k$. 
	Then $E$ is a trivial skew brace with cyclic additive group.
\end{cor}

\begin{proof}
	Using Proposition~\ref{solu1}, we have 
	\[
	E \cong (\cdots((\C_{p_1} \rtimes_{Sb} \C_{p_2}) \rtimes_{Sb} \C_{p_3}) \rtimes_{Sb} \cdots) \rtimes_{Sb} \C_{p_k}.
	\]
	Since $p_2 \nmid (p_1 - 1)$, there cannot exist a nontrivial homomorphism 
	\[
	\C_{p_2} \to \Aut(\C_{p_1}).
	\]
	Hence, the associated actions $\nu$, $\mu$, and $\sigma$ defined in Theorem~\ref{splitsbext} are trivial in this case. Therefore, the skew brace operations defined in \eqref{sb+} and \eqref{sbcirc} coincide, and thus 
	\[
	\C_{p_1} \rtimes_{Sb} \C_{p_2} = \C_{p_1} \times \C_{p_2}
	\]
	is a trivial skew brace of abelian type.  
	
	Next, since $\Aut(\C_{p_1} \times \C_{p_2}) = \Aut(\C_{p_1}) \times \Aut(\C_{p_2})$ (as $p_1$ and $p_2$ are coprime), any nontrivial homomorphism 
	\[
	\C_{p_3} \to \Aut(\C_{p_1}) \times \Aut(\C_{p_2})
	\]
	will surely induce a nontrivial homomorphism 
	\[
	\C_{p_3} \to \Aut(\C_{p_1}) \quad \text{or} \quad \C_{p_3} \to \Aut(\C_{p_2}),
	\]
	which is not possible, since $p_3$  neither divides $p_1 - 1$ nor $p_2 - 1$. Hence, there is no nontrivial homomorphism 
	\[
	\C_{p_3} \to \Aut(\C_{p_1} \times \C_{p_2}).
	\]
	Again using  Theorem~\ref{splitsbext}, we get the  trivial skew brace of abelian type
	\[
	(\C_{p_1} \times \C_{p_2}) \rtimes_{Sb} \C_{p_3} = (\C_{p_1} \times \C_{p_2}) \times \C_{p_3}.
	\]
  Proceeding inductively, we obtain
	\[
	E \cong \C_{p_1} \times \C_{p_2} \times \C_{p_3} \times \cdots \times \C_{p_k},
	\]
	which is a  trivial skew brace  of abelian type. More specifically, $(E, +)$ is a  cyclic group, which completes the proof.
\end{proof}



\begin{prop}
	Let $E$ be a finite skew brace such that $|E / \Ann(E)|$ and $|\Ann(E)|$ are coprime. Then $E \cong \Ann(E) \times H$
	for some sub-skew brace $H$ of $E$ having order $|E / \Ann(E)|$.
\end{prop}

\begin{proof}
	By Theorem~\ref{main7}(3), we have 
	\[
	E \cong \Ann(E) \rtimes_{(\nu, \mu, \sigma)} H,
	\] 
	where $|H| = |E / \Ann(E)|$. It is easy to see that the actions $\nu$, $\mu$, and $\sigma$ become trivial when considering the annihilator extension. The result then follows directly from  \eqref{sb+} and \eqref{sbcirc}.
\end{proof}


Let $E$ be a skew brace having an ideal $I$, which admits a supplement $G$ in $E$.   Then we call $E$ the {\it reduced product} of $I$ and  $G$ if $G$ admits no proper sub-skew brace which supplements $I$ in $E$. The definition specializes to groups if we consider trivial skew braces.  Let us see an example of a reduced product. 
\begin{example}\label{exp1}
	Let $A$ and $I$  be skew braces  as defined in Example \ref{exp0}. Let $G = \langle (0,1) \rangle$. It is easy to see that $G$ is a sub-skew brace of $A$, and we have $I + G = A$. The only proper sub-skew brace $K$ of $G$ is $  \langle (0,2) \rangle$, and we have $I + K = I$. This shows that $A$ is the reduced product of $G$ by $I$.
\end{example}

Let $I$ be an abelian sub-skew brace of $E$ which admits a supplement $G$. Then $J := I \cap G$ is an ideal of $E$ (\cite[Lemma 28]{BEJP24}).  Now $G/J$ is a complement of $I/J$ in $E/J$. 

 We get the following connection between $E$ and $\Lambda_E$ regarding reduced products.

\begin{prop}
Let $E$ be a skew brace having an ideal $I$ and a sub-skew brace $H$. If $\Lambda_E$ is the reduced  product of $\Lambda_I$ and $\Lambda_H$, then $E$ is  the reduced product of $I$ and $H$.
\end{prop}
\begin{proof}
Let $\Lambda_E = \Lambda_I  \cdot  \Lambda_H$ be the reduced product.  Then for any $(a,0) \in \Lambda_E$,  there exist $x, y \in I$ and $h, g \in H$ such that
$$(a,0) = (x,y)(h,g) = (x + \lambda_y(h), y \circ g) = \left((x + \lambda_y(h) - h) + h, y \circ g\right).$$
As $\lambda_y(h) -h \in I$, this implies $a \in I + H$. So $H$ supplements $I$ in $E$.

Contrarily, assume that  there exists a sub-skew brace $K < H$  of $E$ such that $I + K = E$. Then an easy computation show that $\Lambda_E = \Lambda_I \cdot \Lambda_K$, which contradicts the minimality of $\Lambda_H$ as $\Lambda_K < \Lambda_H$. Hence $E$ is the reduced product of $I$ and $H$, completing the proof.
\end{proof}


Let $E$ be a skew brace. The \emph{Frattini sub-skew brace} of $E$, denoted by $\Phi(E)$, is defined as the intersection of all maximal sub-skew braces of $E$, if the intersection exists. Otherwise we define it as   $\Phi(E) = E$.


\begin{prop}\label{prop-reduced}
Let $E = I + H$ be a reduced  product of $I$ and $H$. Then
	\begin{itemize}
		\item[(i)] $I \cap H \leq \Phi(H)$,
		\item[(ii)] $\Phi(E/I) \cong (I + \Phi(H))/I$.
	\end{itemize}
\end{prop}
\begin{proof}
 Let \( C := I \cap H \not\leq \Phi(H) \). Then there exists a maximal sub-skew brace \( M \subsetneq H \) such that \( C \not\leq M \). Clearly, as \( I \) is an ideal of \( E \), the set \( C \) is an ideal of \( H \). Using the maximality of \( M \), we have \( C + M = H \). Hence,
\[
E = I + H = I + (C + M) = (I + C) + M = I + M,
\]
which contradicts the minimality of \( H \), and (i) holds.

Now we prove (ii). Using \( E = I + H \), we have $E/I \cong H / (H \cap I).$ This implies
\[
\Phi(E/I) \cong \Phi(H / (H \cap I)),
\]
where \( \Phi(H / (H \cap I)) \) is the intersection of all maximal sub-skew braces of \( H \) containing \( H \cap I \). But using part (i), \( H \cap I \) is already contained in all maximal sub-skew braces of \( H \). Hence
\[
\Phi(H / (H \cap I)) = \Phi(H) / (H \cap I).
\]
Now it is easy to see that \( H \cap I = \Phi(H) \cap I \), so
\[
\Phi(H) / (H \cap I) = \Phi(H) / (\Phi(H) \cap I) \cong (I + \Phi(H)) / I.
\]
This completes the proof.
\end{proof}

\begin{cor}
Let  $0 \longrightarrow I \xrightarrow{i} E \xrightarrow{\pi} H \longrightarrow 0$
be a brace extension with $E$ finite. Then there exists a sub-skew brace $B$ of $E$ such that $\pi(\Phi(B)) \cong \Phi(H)$.
\end{cor}

\begin{proof}
	If $I \leq \Phi(E)$, then, we have $\Phi(E/I) = \Phi(E)/I$. Since $H \cong E/I$, it follows that $\Phi(H) \cong \Phi(E)/I = \pi(\Phi(E))$, and we can take $B = E$. So, assume that $I \nleq \Phi(E)$. Then there exists a sub-skew brace $B$ of $E$ such that $E$ is the reduced product of $I$ and $B$.  Let $\bar{\pi} : E/I \rightarrow H$ be the isomorphism defined by $\bar{\pi}(\bar{x}) = \pi(x)$ for all $x \in E$. Since $\bar{\pi}$ is an isomorphism, we get
	\[
	\Phi(H) \cong \Phi(E/I) \cong \bar{\pi}(\Phi(E/I)).
	\]
	From assertion (ii) of the preceding result,  we know  that
	\[
	\Phi(E/I) \cong (I + \Phi(B)) / I.
	\]
Hence
	\[
	\Phi(H) \cong \bar{\pi}((I + \Phi(B))/I) = \pi(\Phi(B)),
	\]
	as required, and the proof is complete.
\end{proof}



Let  $0 \longrightarrow I \xrightarrow{i} E \xrightarrow{\pi} H \longrightarrow 0$ be a brace extension. Then $E$ is called a \emph{minimal extension} if $E$ does not admit any proper sub-skew brace $K$ such that $\pi(K) = H$.  Equivalently, $E$ is  said to be minimal with respect to $I$,  if $I$ does not admit a proper supplement in $E$. We remark that a minimal brace extension is on the extreme end of the split brace extension. We first produce two examples of  minimal extensions.

\begin{example}
	Let $A$ be a skew brace  as in Example \ref{exp0}. It is easy to see that $\Soc(A) = \langle (0,2) \rangle^+$. GAP computations show that there is no proper sub-skew brace $K$ of $A$ such that $K + \Soc(A) = A$. Hence, $A$ is a minimal extension with respect to $\Soc(A)$.
\end{example}	

\begin{example}
Let $A$ be a  brace with additive group $\mathbb{Z}_8$. Define  `$\circ$' on $\mathbb{Z}_8$ as follows:
\[
x \circ y = x + 5^{x}y \pmod{8}.
\]
It turns out that $(A,\circ) \cong \mathbb{Z}_8$ as classified in \cite{Bach15}.  
It is easy to see that $\Soc(A)=\langle 2 \rangle$. Using $\lambda_x(y)=5^x y \pmod{8}$, we note that
\[
\text{if $x$ is odd, then } 5^{x} \equiv 5 \pmod{8}
\mbox{ and }
\text{if $x$ is even, then } 5^{x} \equiv 1 \pmod{8}.
\]
Therefore, every subgroup of $\mathbb{Z}_8$ is an ideal of $A$.  
Since $\Soc(A)$ is a maximal ideal of $A$, this shows that $A$ is a minimal extension of $\Soc(A)$.
\end{example}

The following two results are other  easy consequences of Proposition \ref{prop-reduced} on simple extensions.
\begin{cor}
A brace extension $E$ of a skew brace $H$ by an ideal $I$ is minimal if and only if $I \leq \Phi(E)$.
\end{cor}

\begin{cor}
	Let $E$ be a brace  extension of $H$ by $I$. If $\Lambda_E$ is the minimal extension of $\Lambda_H$ by $\Lambda_I$, then $E$ is also a minimal extension of $H$ by $I$.
\end{cor}


The converse of the above corollary need not be true.  
Let $p$ be any prime. Since every skew brace of order $p^2$ is a minimal extension of $\mathbb{Z}_p$ by $\mathbb{Z}_p$, consider 
$A := \texttt{SmallSkewbrace}(4,2)$ and let $I := \mathbb{Z}_2$.  GAP computations show that 
\[
\Lambda_A \cong \mathbb{Z}_2 \times D_4,
\] 
where $D_4$ is the dihedral group of order $8$, which is \emph{not} a minimal extension of 
\[
\Lambda_I = \mathbb{Z}_2 \times \mathbb{Z}_2.
\]


\begin{prop}
	Let $E$ be a minimal extension of a skew brace $H$ with respect to  $I$ and let $J \leq I$ be such that  $J$ is an ideal of $E$. Then $E/J$ is a minimal extension of $H$ with respect to $I/J$. 
\end{prop}

\begin{proof}
Suppose that \( E/J \) is not a minimal extension of \( H \) by \( I/J \). Then there exists a sub-skew brace \( K \subsetneq E \) containing \( J \) such that \( K/J + I/J = E/J \). Hence, \( E = K + J \), contradicting the minimality of \( E \).
\end{proof}

There do exist brace extensions which admit only extreme states, that is, minimal or split. 
\begin{prop}
Let $I$ be a simple $H$-module, where $H$ is skew brace. Then the corresponding extension $E$, say,  of $H$ by $I$  either splits  or is minimal. 
\end{prop}
\begin{proof}
Suppose that \( E \) is not a minimal extension. Then $I$ admits a proper supplement $G$, say, in $E$.  Since $I$ is an ideal of $E$, $G \cap I$ is an ideal of $G$.  We claim that \( G \cap I \) is an ideal of \( E \).  For any  \( h \in E \),  we know that \( h = g \circ z \) for some \( g \in G \), \( z \in I \). Then, for any \( x \in G \cap I \), we get
	\[
	h \circ x \circ h^{-1} = g \circ z \circ x \circ z^{-1} \circ g^{-1} = g \circ x \circ g^{-1}
	\]
belongs to $G$ as well as $I$. Thus \( G \cap I \) is normal in \( (E, \circ) \). Similarly,  \( G \cap I \) is a normal subgroup of \( (E, +) \). Also
	\[
	\lambda_{g \circ z}(x) = \lambda_g(x)
	\]
belongs to $G$ as well as $I$. Hence \( \lambda_h(x) \in G \cap I \), and therefore the claim follows.   Thus, \( G \cap I \) is an \( H \)-submodule of \( I \), which, by the given hypothesis,  is possible only when $G \cap I = \{0\}$ as, $G$ being proper supplement, $G \cap I \ne I$. Hence $G$ is a complement of $I$ in $E$, which  implies that the brace extension \( E \) is a split extension, and the proof is complete.
	\end{proof}

 Concrete case is
\begin{cor}
Let $E$ be an extension of a skew brace $H$ by a cyclic group $I$ of prime order (as a simple skew brace). Then, either $E$ is a split extension or a minimal extension.
\end{cor}


%%%%%
Before proceeding further, we define a generating set of a skew brace. We follow the constructions done in \cite{CFT25, MT23}.  Let $A$ be a skew brace and $X \subseteq  A$.  Denote by $A(X)$ the smallest sub-skew brace of $A$ containing  $X$.  We say that the \emph{ set $X$  generates  $A$}  if  $A = A(X)$ and we write $A = \gen{X}$.  We now investigate the general form of the elements of a sub-skew brace of a skew brace $A$ generated by a finite subset  
$X$ of $A$.

Let  $M$ denote the free monoid generated by the set $Y := \{+, \bar{+}, \circ, \bar{\circ}, \delta_1,  \delta_2, \chi_1, \chi_2\}$. 
Any $m \in M$ has the form $m = \epsilon_1  \epsilon_2 \cdots  \epsilon_s$, where  $\epsilon_i \in Y$ and $s$ is a non-negative integer. 
We define \emph{the degree of $m$} to be the number of $\epsilon$'s which belongs to $\{+, \bar{+}, \circ, \bar{\circ}\}$, and denote it by $\deg(m)$. Given a skew brace $A$, for $a, b \in A$, we set 
$$+(a, b) = a+b, ~ \bar{+}(a, b) = b+a, ~ \circ(a, b) = a \circ b, ~ \bar{\circ}(a, b) = b \circ a$$
and 
$$\delta_1(a, b) = (-a, b), ~ \delta_2(a, b) = (a, -b), ~ \chi_1(a, b) = (a^{-1}, b), ~ \chi_2(a, b) = (a, b^{-1}).$$
Let $m =  \epsilon_1 \cdots \epsilon_{i-1} \epsilon_{1}' \epsilon_{i} \cdots  \epsilon_{j-1} \epsilon_{r}' \epsilon_{j} \cdots  \epsilon_t \in M$ have degree $t$, where $\epsilon_1, \ldots, \epsilon_t \in \{ +, \bar{+}, \circ, \bar{\circ}\}$ and  $\epsilon_1', \ldots,  \epsilon_{r}'  \in  
\{\delta_1, \delta_2,  \chi_1, \chi_2\}$.
For the given  elements  $a_1, a_2, \ldots, a_{t+1} \in A$, we define
\begin{equation}
m(a_1, \ldots, a_{t+1}) = \epsilon_1(a_{1}, \epsilon_2(a_2, \ldots \epsilon_{j-1}( \epsilon_{r}'(a_{j-1}, (\epsilon_j(a_j, \cdots \epsilon_t(a_t, a_{t+1}) \cdots )))) \cdots)).
\end{equation}


Let $A$ be a skew brace generated  by a finite subset $X := \{x_1, x_2, \ldots, x_r\}$ of $A$.  Let   $-X := \{-x \mid x \in X\}$ and $X^{-1} := \{x^{-1} \mid x \in X\}$. Set $X_1 := X \cup -X \cup X^{-1}$ and  define 
\[
K_{1} := X_1  \cup \{x+y,\; x \circ y \mid  x, y \in X_1\}. 
\]
Now set  $X_2 := K_1 \cup -K_1 \cup K_1^{-1}$ and define
$$K_2 := X_2  \cup \{x+y,\; x \circ y \mid  x, y \in X_2\}.$$
Next, for $n \ge 3$, inductively define
\begin{eqnarray*}
X_n &:=&  K_{n-1} \cup -K_{n-1} \cup K_{n-1}^{-1}\\
 K_n &:=& X_n \cup   \{x+y,\; x \circ y \mid  x, y \in X_n\}.
 \end{eqnarray*}
 Finally set
\[
\mathcal{W}(X):=\bigcup_{n\ge1} K_n.
\]
Let $x \in K_n$. Then there exist elements $y_1, y_2 \in K_{n-1}$ such that 
 $x = \epsilon \epsilon'(y_1, y_2)$, where $\epsilon \in   \{+, \bar{+}, \circ, \bar{\circ}\}$, $\epsilon' \in \{1_M, \delta_1, \delta_2,  \chi_1, \chi_2\}$,   $1_M$ being the  identity element of $M$.


With this setting, we have the following interesting result.

\begin{prop}\label{helping}
The following statements hold true:
\begin{enumerate}
\item[(i)] The set  \(\mathcal W(X)\) is a sub-skew-brace of \(A\) generated by $X$. 
\item[(ii)]  For every \(0 \ne x \in \mathcal{W}(X)\), there exist  an  integer \(n \geq 0\),  elements
$$\epsilon_{0\blacktriangle}, \epsilon_{1\blacktriangle}, \ldots, \epsilon_{(n-1)\blacktriangle} \in \{+, \bar{+}, \circ, \bar{\circ}\}$$
and 
$$\epsilon'_{0\blacktriangle}, \epsilon'_{1\blacktriangle}, \ldots, \epsilon'_{(n-1)\blacktriangle} \in \{1_M, \delta_1, \delta_2,  \chi_1, \chi_2\}$$
such that 
$$x = \epsilon_{(n-1) \blacktriangle} \epsilon'_{(n-1) \blacktriangle}\big((\cdots \epsilon_{1 \blacktriangle} \epsilon'_{1 \blacktriangle}(\epsilon_{0\blacktriangle} \epsilon'_{0 \blacktriangle}(\bullet, \bullet), \epsilon_{0\blacktriangle} \epsilon'_{0 \blacktriangle}(\bullet, \bullet))), (\cdots \epsilon_{1 \blacktriangle} \epsilon'_{1 \blacktriangle}(\epsilon_{0\blacktriangle} \epsilon'_{0 \blacktriangle}(\bullet, \bullet), \epsilon_{0\blacktriangle} \epsilon'_{0 \blacktriangle}(\bullet, \bullet)))\big),$$
where $\bullet \in X \cup \{0\}$ and may take different values at different positions. Also, for any brace homomorphism $\phi : A \to B$, the image of $x$ under $\phi$, denoted by $\phi(x)$, is given by 
$$\epsilon_{(n-1) \blacktriangle} \epsilon'_{(n-1) \blacktriangle}\big((\cdots \epsilon_{1 \blacktriangle} \epsilon'_{1 \blacktriangle}(\epsilon_{0\blacktriangle} \epsilon'_{0 \blacktriangle}(\bullet', \bullet'), \epsilon_{0\blacktriangle} \epsilon'_{0 \blacktriangle}(\bullet', \bullet'))), (\cdots \epsilon_{1 \blacktriangle} \epsilon'_{1 \blacktriangle}(\epsilon_{0\blacktriangle} \epsilon'_{0 \blacktriangle}(\bullet', \bullet'), \epsilon_{0\blacktriangle} \epsilon'_{0 \blacktriangle}(\bullet', \bullet')))\big),$$
where $\bullet' = \phi(\bullet) \in B$.
\end{enumerate}
\end{prop}

\begin{proof}
 	 Obviously, \(X =\{x_{1},\dots,x_{r}\} \subseteq  \mathcal{W}(X)\).  For any two elements  \(x,y\in\mathcal{W}(X)\),  there exists some integer  \(n\geq 0\) such that \(x,y \in K_{n}\). It follows that 
	\[
	-x,\quad x^{-1},\quad x+y,\quad x\circ y \in K_{n+1}\subseteq \mathcal{W}(X).
	\]
	Thus \(\mathcal{W}(X)\) is closed under additive inverses, multiplicative inverses, and both operations. Therefore, \(\mathcal{W}(X)\) is a sub-skew brace of \(B\) containing $X$.  It is easy to see by induction that \(K_{n}\subseteq B(X)\) for all \(n \geq 1\). Consequently,
	\[
	\mathcal{W}(X)=\bigcup_{n\geq 0}K_{n}\subseteq B(X),
	\]
	which proves that  $ \mathcal{W}(X)=B(X)$.
	
The second assertion follows from the definition of  $\mathcal{W}(X)$, definition of a brace homomorphism  and applying an easy induction.  Indeed, for the induction base, we see that $x = +(0, x)$, $-x = + \delta_2(0, x)$ and $x^{-1} = \circ \chi_2(0, x)$ for the elements of $X_1$.  Note that these expressions are not unique. So the base case when $0 \ne x \in K_1$ is clear. 
\end{proof}

We remark that the element $0 \ne x \in \mathcal{W}(X)$ as in the preceding proposition can be expressed as
$$x = \epsilon_{(n-1) \blacktriangle} \epsilon'_{(n-1) \blacktriangle} \epsilon_{(n-2) \blacktriangle} \epsilon'_{(n-2)\blacktriangle} \cdots \epsilon_{1 \blacktriangle} \epsilon'_{1 \blacktriangle} \epsilon_{0\blacktriangle} \epsilon'_{0 \blacktriangle}\big(a_{n-1} a_{n-2} \cdots a_1 a_{01}a_{02}\big),$$
where $a_{n-1} \in K_{n-1}, a_{n-2} \in K_{n-2} , \ldots, a_1 \in K_1, a_{01}, a_{02}  \in X$.  For the convenience of the notation, let us denote this expression in word map form as follows: Let $0 \ne x \in \mathcal{W}(X)$ as in the preceding proposition. Then we can write $x$ as an image of a word $w: X^m \to  \mathcal{W}(X)$, that is, $x = w(x_{i_1}, \ldots, x_{i_m})$, where $m \le 2^n$ and $x_{i_j} \in X$. More precisely, $w$ is the expression
$$\epsilon_{(n-1) \blacktriangle} \epsilon'_{(n-1) \blacktriangle}\big((\cdots \epsilon_{1 \blacktriangle} \epsilon'_{1 \blacktriangle}(\epsilon_{0\blacktriangle} \epsilon'_{0 \blacktriangle}(\bullet, \bullet), \epsilon_{0\blacktriangle} \epsilon'_{0 \blacktriangle}(\bullet, \bullet))), (\cdots \epsilon_{1 \blacktriangle} \epsilon'_{1 \blacktriangle}(\epsilon_{0\blacktriangle} \epsilon'_{0 \blacktriangle}(\bullet, \bullet), \epsilon_{0\blacktriangle} \epsilon'_{0 \blacktriangle}(\bullet, \bullet)))\big),$$
where $\bullet$'s are indeterminants $t_1, t_2, \ldots, t_m$.

\begin{thm}
Let $0 \longrightarrow I \xrightarrow{i} E \xrightarrow{\pi} H \longrightarrow 0$ be a minimal extension of a skew brace $H$ by $I$.  If $H$ is an $r$-generated skew brace, then $E$ is also an $r$-generated skew brace.
\end{thm}

\begin{proof}
Let $H$ be generated by $X=\{x_1,x_2, \ldots, x_r\}$, and let  $s \colon H \to E$ be a set-theoretic section of $\pi$.  
Let $T$ be the sub-skew brace of $E$ generated by  $s(X)= \{s(x_1), s(x_2), \ldots, s(x_r)\}$.  
For any arbitrarily fixed element $x \in H$, by Proposition~\ref{helping}(1), there exists a word $w$  and an integer $m$ (both depending on $x$) such that 
\[
x = w(x_{i_1},\ldots,x_{i_m}) 
\]
for some $x_{i_j} \in X$, $1 \leq j \leq m$.  
Applying Proposition~\ref{helping}(2), we obtain
\begin{align}\label{conversion}
s(w(x_{1_i}, \ldots, x_{i_m})) - w(s(x_{i_1}),\ldots,s(x_{i_m})) \in I.
\end{align}
Indeed,
\begin{eqnarray*}
\pi\bigl(s(w(x_{1_i}, \ldots, x_{i_m})) - w(s(x_{i_1},\ldots,s(x_{i_m}))\bigr)
&=&   w(x_{i_1},\ldots,x_{i_m}) - w(\pi(s(x_{i_1})), \ldots, \pi(s(x_{i_m})))\\
&=&  w(x_{i_1},\ldots,x_{i_m}) - w(x_{i_1},\ldots,x_{i_m}) = 0 .
\end{eqnarray*}

Now let $a \in E$ be an arbitrary element.   We know that $a = s(x) + y$ for some $x \in H$ and $y \in I$.  
Writing $x = w(x_{i_1},\ldots,x_{i_m}) $ and using \eqref{conversion}, we get
\[
a = s(w(x_{i_1},\ldots,x_{i_m}) ) + y 
= w(s(x_{i_1}), \ldots,s(x_{i_m})) + z + y,
\]
for some $z \in I$.  
This shows that $a \in T+I$.  
Since $E$ is a minimal extension, we must have $T = E$.  
Hence, $E$ is an $r$-generated skew brace.
\end{proof}


A kind of converse of the preceding theorem also holds true.

\begin{thm}\label{confrat}
Let $0 \longrightarrow I \xrightarrow{i} E \xrightarrow{\pi} H \longrightarrow 0$ be an extension of a skew brace $H$ by $I$.   If $H$ and $E$ are  both  minimally generated by $r$ elements, then  $I \subseteq \Phi(E)$, and therefore  $E$ is a minimal extension of $H$ by $I$.
\end{thm}


\begin{proof}
	Without loss of generality, assume that $\Phi(E) \neq E$. Suppose, for the sake of contradiction, 
	that $I \not\subseteq \Phi(E)$. Then there exists an element $x \in I \setminus \Phi(E)$. 
	Thus, we can extend $\{x\}$ to a generating set $S = \{ g_1, \ldots, g_{r-1}, x\}$ of  $E$, 
	but no proper subset of $S$ generates $E$.  It is sufficient to prove that  $\pi(S) = \{ h_1 :=\pi(g_1), \ldots, h_{r-1} := \pi(g_{r-1}), \pi(x)=0\}$ generates $H$, as it will contradict the given supposition that $H$ is minimally generated by $r$ elements.
	
	 For any  $h \in H$, the element   $s(h) \in E$  can be represented as 
	\[
	s(h) = w(g_{i_1}, \ldots, g_{i_m})
	\] 
	for some word $w : S^m \to E$ as explained above, where $g_{i_j} \in S$.
	Applying the projection map $\pi: E \to H$, we get
	\begin{eqnarray*}
	h &=& \pi(s(h)) =\pi\big( w(g_{i_1}, \ldots, g_{i_m})  \big)\\
	&=&  w(\pi(g_{i_1}), \ldots, \pi(g_{i_m})) = w(h_{i_1}, \ldots, h_{i_m}),
	\end{eqnarray*}
where $h_{i_j} \in \pi(S)$, $1 \le j \le m$. This shows  that $H$ is generated by $\pi(S)$, and the proof is complete.
\end{proof}





\begin{prop}
	Let $E$ be a minimal extension of a skew brace $H$ by $I$.   If the rank of $(H,+)$ (respectively of $(H,\circ)$) is $r$,  then the rank of $(E,+)$ (respectively of $(E,\circ)$) is at most $r$.
	
	\end{prop}
\begin{proof}
Let $(H,+)$ be generated by $\{h_1,h_2,\ldots,h_r\}$, and let 
$s\colon H \to E$ be any set-theoretic section.  
Let $T$ be the subgroup of $(E, +)$ generated by 
$\{s(h_1),s(h_2),\ldots,s(h_r)\}$.  
It is easy to see that  
\[
s\, \Bigl(\sum_{j=1}^{t}k_j h_{i_j}\Bigr)-
\sum_{j=1}^{t}k_j s(h_{i_j})\in I 
\quad\text{for all integers }k_j \text{ and } h_{i_j} \in \{h_1,h_2,\ldots,h_r\}.
\]
Thus there exists some $z\in I$, depending on the $h_{i_j}$, and $k_j$, such that  
\[
s\, \Bigl(\sum_{i=j}^{t}k_i h_{i_j}\Bigr)
= \sum_{j=1}^{t}k_j s(h_{i_j}) + z.
\]
 We know that any $x \in E$ can be written as $x=s(h)+y$ for some $h \in H$ and $y \in I$.  
We can write $h=\sum_{j=1}^{t}k_i h_{i_j}$ for some integers $k_j$ and $h_{i_j} \in \{h_1,h_2, \ldots, h_r\}$.  
This gives 
\[
x=s\, \Bigl(\sum_{i=1}^{m}k_i h_i\Bigr)+y
= \sum_{i=1}^{m}k_i s(h_i)+z+y,
\]
which shows that $x \in T+I$.   Since $E$ is a minimal extension, we have $T=E$.  This shows that the rank of $(E,+)$ is at most $m$. Similarly we get the assertion for $(E, \circ)$, and the proof is complete.
\end{proof}

A skew brace $(A, +, \circ)$ is said to be {\it cyclic} ({\it cocyclic}) if $(A, +)$ ($(A, \circ)$) is a cyclic group.
\begin{cor}
Let $E$ be a minimal extension of a skew brace $H$ by $I$. If $H$ is a cyclic (respectively, cocyclic) skew brace, then $E$ is also a cyclic (respectively, cocyclic) skew brace.
\end{cor}



\begin{cor}
	Let $E$ be a finite supersoluble  non-zero cyclic skew brace with a chief series
	\[
	\{0\} = I_0 \leq I_1 \leq \cdots \leq I_k = E.
	\] 
 Then
	\[
	I_i \leq \Phi(I_{i+1}) \quad \text{for all } 0 \leq i < k.
	\] 
	In particular, $I_{k-1} = \Phi(E)$.
\end{cor}



\begin{proof}
	Since $(E, +)$ is cyclic, the additive groups of $I_i$ are also cyclic for all $0 \leq i \leq k$. 
	Hence, each $I_{i+1}$ is a $1$-generated skew brace such that $I_{i+1}/I_i$ is also $1$-generated. 
	Using Theorem~\ref{confrat}, we  get
	\[
	I_i \leq \Phi(I_{i+1}).
	\] 
As $E$ is a non-zero finite  skew brace, it follows that $\Phi(E) \neq E$.  Moreover, since $E/I_{k-1}$ has no non-zero sub-skew brace,  $I_{k-1}$ is a maximal sub-skew brace of $E$. Now, $(E,+)$ being cyclic, $I_{k-1}$ is the unique  maximal sub-skew brace of $E$. Hence $I_{k-1} = \Phi(E)$, and the proof is complete.	 
\end{proof}

 For any two skew  braces $E_1$ and $E_2$, let $E:=E_1 \oplus E_2$ be generated by a subset $X$ of $E$. Then it is not difficult to see that $E_1 = \gen{\pi_1(X)}$ and $E_2 = \gen{\pi_2(X)}$, where $\pi_i :  E \to E_i$, $i = 1, 2$, are the natural projections of skew braces.  Now we prove

\begin{thm}
Let $E = E_1 \oplus E_2$ be a finite skew brace. Then $\Phi(E) = \Phi(E_1) \oplus \Phi(E_2)$. 
\end{thm}
\begin{proof}
For every maximal sub-skew brace $M_1$ of $E_1$, $M_1 \oplus E_2$ is a maximal sub-skew brace of $E$. Thus  $\Phi(E) \le  \Phi(E_1) \oplus E_2$. Similarly $\Phi(E) \le  E_1 \oplus \Phi(E_2)$. Thus  $\Phi(E) \le  \Phi(E_1) \oplus \Phi(E_2)$. Next we prove that $\Phi(E_1) \oplus \Phi(E_2) \leq \Phi(E)$.  

Let $ \Phi(E_1) \neq E_1$  and   $\Phi(E_2) \neq E_2$. We show that every element $(a,b) \in \Phi(E_1) \oplus \Phi(E_2)$ is a non-generator in $E_1 \oplus E_2$.  Let $(a,b) \in \Phi(E_1) \oplus \Phi(E_2)$ and suppose that $X$ is a subset of $E_1 \oplus E_2$ such that $\langle X \cup \{(a,b)\} \rangle = E_1 \oplus E_2$. It suffices to prove $E_1 \oplus E_2 = \langle X \rangle $.  Consider the projection maps $\pi_{1}: E_1 \oplus E_2 \to E_1$ and $\pi_{2}: E_1 \oplus E_2 \to E_2$. Since $\langle X \cup \{(a,b)\} \rangle = E_1 \oplus E_2$, we have
\begin{align*}
	E_1 &= \pi_{1}(\langle X \cup \{(a,b)\} \rangle) = \langle \pi_{1}(X) \cup \{a\} \rangle, \\
	E_2 &= \pi_{2}(\langle X \cup \{(a,b)\} \rangle) = \langle \pi_{2}(X) \cup \{b\} \rangle.
\end{align*}

Since $a \in \Phi(E_1)$, we have $E_1 = \langle \pi_{1}(X)  \rangle$. Similarly, since $b \in \Phi(E_2)$, we get  $E_2 = \langle \pi_{2}(X) \rangle$.  Now,  for any $(x,y) \in E_1 \oplus E_2$, we have $x \in E_1 = \langle \pi_{1}(X) \rangle$ and $y \in E_2 = \langle \pi_{2}(X) \rangle$. Then there exist elements $x_{i_1}, \ldots, x_{i_k},  x_{j_1}, \ldots, x_{j_l} \in X$ (with possible repeatitions) and words $w_1, w_2$ such that
\begin{align*}
	x &= w_1\big(\pi_{1}(x_{i_1}), \pi_{1}(x_{i_2}), \ldots, \pi_{1}(x_{i_k})\big), \\
	y &= w_2\big(\pi_{2}(x_{j_1}), \pi_{2}(t_{j_2}), \ldots, \pi_{2}(x_{j_l})\big),\\
0 &= w_1\big(\pi_{2}(x_{i_1}), \pi_{2}(x_{i_2}), \ldots, \pi_{2}(x_{i_k})\big),\\
0 &= w_2\big(\pi_{1}(x_{j_1}), \pi_{1}(t_{j_2}), \ldots, \pi_{1}(x_{j_l})\big).
\end{align*}
Note that, in $E$,  $(x, 0)$ is a word in  $x_{i_1}, \ldots, x_{i_k}$  and $(0, y)$ is a word in $x_{j_1}, \ldots, x_{j_l}$.

Therefore, the element $(x,y)$ can be expressed as
\[
\big(w_1(\pi_{1}(x_{i_1}), \ldots, \pi_{1}(x_{i_k})) + w_2 (\pi_{1}(x_{j_1}), \ldots, \pi_{1}(x_{j_l})), \,
w_1 (\pi_{2}(x_{i_1}), \ldots, \pi_{2}(x_{i_k})) + w_2(\pi_{2}(x_{j_1}), \ldots, \pi_{2}(x_{j_l})\big),
\]
which, considering the fact that $z = \big(\pi_1(z), \pi_2(z)\big)$ for all $z \in X$, is an image of the  word $\epsilon w_1w_2$ on the elements of $X$, where $\epsilon = +$. Thus $(x,y) \in  E(X)$ and therefore $E = \gen{X}$. The proof is complete.
\end{proof}

It is well know that for a finite $p$-group $G$, $\Phi(G) = \gamma_2(G)G^p$, where $G^p$ denotes the subgroup of $G$ generated by the set $\{g^p \mid g \in G\}$. It is intriguing to obtain the structure of the Frattini sub-skew brace of a skew brace of prime power order.

  Let $E$ be a cyclic skew brace of order not equal to any prime. Then $E$ admits a non-zero proper sub-skew brace (\cite[Theorem A]{BEJP24}), and therefore $E$ admits a unique maximal sub-skew brace equal to $\Phi(E)$.  So if $E$ is such a cyclic (cocyclic) skew brace, then $\Phi(E) \le \Phi(E, +)$ ($\Phi(E) \le \Phi(E, \circ)$). But, in general, there does not seem any clear connection. It will be interesting to explore relations among  $\Phi(E)$, $\Phi(E, \circ)$ and $\Phi(E, +)$.
  
  


%%%%%%%%%%%%%%%%%%%%%%%%%%%%%%%%%%%%%%%%%%%%%%%%%%%%%%
\section{Reductions for split extension}

	
In this section we study several reduction results on split extensions of finite skew braces by abelian groups (viewed as a trivial  skew brace of abelian type). These results are brace analogs of group theory results from \cite{Hill72}. The following basic result is well known.
\begin{lemma}[Dedekind's modular law]\label{Distribute}
Let $G$ be a group, and let $X$, $Y$ and $T$ be subgroups of $G$ such that $T$ is contained in $X$. Then
\[
(X \cap Y)T = XT \cap YT.
\]
\end{lemma}

We now prove a reduction result for split brace extensions.
\begin{thm}\label{thm1}
	 Let $E$ be a brace  admitting an ideal  $I$ such that $I = I_1 \times \cdots \times I_r$ with $I_j \trianglelefteq E$ for $1 \leq j \leq r$. Let $I_k^* = \prod_{j \neq k} I_j$. Then the following  hold true:
\begin{itemize}
	\item[(i)] $E$ splits over $I$ if and only if $E/I_k^*$ splits over $I/I_k^*$ for all $k = 1, \ldots, r$.
	\item[(ii)]  $E$ splits over $I$ if and only if, for any $k \neq m$, the quotients $E/I_k$ and $E/I_m$ split over $I/I_k$ and $I/I_m$, respectively.
\end{itemize}
\end{thm}
\begin{proof}
We first prove (i).   Let $E$  split over $I$. Then $E = G + I$ for some sub-skew brace $G$ of $E$ such that $G \cap I = \{0\}$. Then we get
\[
E / I_k^* =( G+I) / I_k^* \quad \text{for } 1 \leq k \leq r,
\]
which gives
\[
E / I_k^*   = ((G + I_k^*) / I_k^*) + (I / I_k^*).
\]
Also note that
\[
((G + I_k^*) / I_k^*) \cap (I / I_k^*) = ((G \cap I) + I_k^*) / I_k^* = I_k^* / I_k^* = \{\bar{0}\}.
\]
Hence, $E / I_k^*$  splits over $I / I_k^*$ with complement  $(G + I_k^*) / I_k^*$ for all $k = 1, \ldots, r$.

Conversely, assume  that $E / I_k^*$ splits over $I / I_k^*$ for each $k = 1, \ldots, r$. Then there exists a sub-skew brace $G_k$ of $E$ containing $I_k^*$ such that 
\[
G_k/ I_k^* +  I / I_k^* = E / I_k^*
\quad \text{and} \quad
G_k \cap I = I_k^*
\]
for each $k$.  Let $G = \bigcap_{k=1}^r G_k$, which is a sub-skew brace of $E$. Then
\[
G + I = \left( \bigcap_{k=1}^r G_k \right) + (I_1 + I_2 + \cdots + I_r)=\left( \left( \bigcap_{k=1}^r G_k \right) + I_1\right)+ (I_2 +  \cdots + I_r).
\]
Using Lemma~\ref{Distribute} and the fact that $G_k + I_k = E$ for $1 \le k \le r$, we get
\[
G + I = \bigcap_{k=1}^r (G_k  + I_k) = E.
\]
Since $G \cap I = \bigcap_{k=1}^r (G_k \cap I) = \bigcap_{k=1}^r I_k^* = \{0\}$,  $E$ splits over $I$.

We now prove (ii). Without loss of generality, let \( k = 1 \) and \( m = 2 \). Suppose that \( E \) splits over \( I \), and let \( G \) be a complement of \( I \) in \( E \). Define
\[
\overline{G}_1 := GI_1 / I_1 \quad \text{and} \quad \overline{I}_1 := I / I_1.
\]
We have \( GI_1 \cap I = I_1 \), so \( \overline{G}_1 \cap \overline{I}_1 = \{I_1/I_1\} \). Hence, \( \overline{G}_1 \) is a complement of \( \overline{I}_1 \) in \( E / I_1 \), and thus \( E / I_1 \) splits over \( I / I_1 \). A similar argument shows the result for \( I_2 \).

Conversely, suppose that \( E/I_1 \) and \( E/I_2 \) split over \( I/I_1 \) and \( I/I_2 \), respectively. Let \( \overline{G}_1 \) and \( \overline{G}_2 \) be complements of \( I/I_1 \) and \( I/I_2 \) in \( E/I_1 \) and \( E/I_2 \), respectively. As we know, \( \overline{G}_1 = G_1/I_1 \) and \( \overline{G}_2 = G_2/I_2 \) for some skew braces \( G_1 \) and \( G_2 \) of \( E \) containing \( I_1 \) and \( I_2 \), respectively. Since \( \overline{G}_1 \) and \( \overline{G}_2 \) are complements of \( I/I_1 \) and \( I/I_2 \), it follows that \( G_1 \cap I = I_1 \) and \( G_2 \cap I = I_2 \). Let \( G_3 = G_1 + (I_3 + I_4 + \cdots + I_r) \), so that \( G_3 + I_2 = E \) and $G_3 \cap I_2=\{0\}$ . Define \( G = G_3 \cap G_2 \). Then we have
\[
G + I = (G_3 \cap G_2) + I = (G_3 \cap G_2) + (I_1 + I_2 + \cdots + I_r).
\]
Using Lemma~\ref{Distribute} repeatedly, we get
\begin{align*}
G + I & = (G_3 + I_1 \cap G_2 + I_1) + (I_2 + \cdots + I_r) \\
&\quad  \vdots\\
&=  (G_3 + (I_1 + I_2 +\cdots +I_r) \cap  (G_2 + (I_1 + I_2 + \cdots  + I_r))\\
& = E \cap E = E.
\end{align*}
Moreover,
\[
G \cap I = G_3 \cap G_2 \cap I = G_3 \cap I_2 = \{0\}.
\]
This shows that \( G \) is a complement of \( I \), and hence \( E \)  splits  over \( I \). The proof is complete.
\end{proof}

Let $I$ be a finite skew brace such that both $(I, +)$ and $(I, \circ)$ are nilpotent groups. Then it follows from \cite[Corollary 4.3]{CSV18} that Sylow subgroups of  $(I, +)$ are ideals of $I$. Furthermore, if $P_i$, $1 \le i \le k$, are all distinct Sylow subgroups of $(I, +)$, then $I = P_1 \oplus P_2 \oplus \cdots \oplus P_k$. With this setting, the preceding  theorem gives
\begin{cor}
Let \( E \) be a brace extension of \( H \) by \( I \), where $I$ is an in the preceding para.   Define \( P_i^* = \oplus_{j \neq i} P_j \) for each \( i \). Then the following hold true:
	\begin{itemize}
		\item[(i)] \( E \) splits over \( I \) if and only if \( E / P_k^* \) splits over \( I / P_k^* \) for all \( k = 1, \ldots, k \).
		\item[(ii)] \( E \) splits over \( I \) if and only if, for any \( k \neq m \), the quotients \( E / P_k \) and \( E / P_m \) split over \( I / P_k \) and \( I / P_m \), respectively.
	\end{itemize}
\end{cor}
\begin{proof}
Since each \( P_i \) is a characteristic subgroup of \( (I, +) \), as observed above,  each of these is an ideal of  \( E \). The assertions are now  immediate from Theorem~\ref{thm1}.
\end{proof}


\begin{remark}
Let $H$ be a skew brace. If \(M\) is an \(H\)-submodule of an $H$-module  \(I\), then \(M \times M\) is a \(\Lambda_H\)-submodule of \(I \times I\). 
\end{remark}

For an $H$-module $I$,  an  \emph{$H$-composition series} for $I$ is a series of submodules of $I$
	\[
	I = I_0 \trianglerighteq I_1 \trianglerighteq \cdots \trianglerighteq I_k = \{e\}
	\]
	such that each quotient $I_i / I_{i+1}$ is a simple $H$-module.


\begin{prop}
	Let $H$ be a skew brace and   $I$ be a finite  $H$-module.  Then there exists an $H$-composition series for $I$, and any two such composition series have the same length. Moreover, their composition factors are the same up to isomorphism and rearrangement.
\end{prop}

\begin{proof}
The existence follows from the finiteness of $I$: if $I$ is not a simple $H$-module, then by the finiteness of $I$, there exists a maximal $H$-submodule of $I$. Repeating this process, we obtain an $H$-composition series for $I$. The second part also follows from a standard group-theoretic argument.
\end{proof}

\begin{cor}\label{refine}
		Let $H$ be a skew brace and   $I$ be a finite  $H$-module. Then any series of $H$-modules can be  refined to an $H$-composition series.
\end{cor}
	


\begin{prop}
Let $E$ be a brace extension of a trivial skew brace  $H$ by a finite abelian $p$-group $I$ and 
\[
I = I_0 \rhd I_1 \rhd \cdots \rhd I_r = \{0\}
\]
be an $H$-composition series for $I$ such that $I_i \trianglelefteq E$ for all $i$. Then each factor $I_i / I_{i+1}$ ($1 \leq i < r$) is an elementary-abelian $p$-group and a simple $H$-module. Furthermore,  if $\sigma$ is the trivial action,  then each $I_i / I_{i+1}$ is a simple $\Lambda_H$-module via the action $\phi^{(\nu, \mu)}$. 
\end{prop}

\begin{proof}
Let $|I| = p^n$. Define $I^{p^{i}} := \gen{z^{p^{i}} \mid z \in I}$. It is well known that $I^{p^{i}} $ is a characteristic subgroup of $I$ for all $i \ge 1$ and $I^{p^{i}}/I^{p^{i+1}}$ is elementary abelian $p$-group. Thus $I^{p^{i}}$ is an ideal of $E$, and therefore an $H$-module. Since $I$ is finite, there exists $k \ge 1$ such that $I^{p^{k}} = \{0\}$. We get a series 
$$I \rhd I^{p^{1}}  \rhd \cdots  \rhd I^{p^{k}} = \{0\}$$
 of $H$-submodules  of $I$.   Using Corollary~\ref{refine}, this series can be refined to an $H$-composition series 
$$ I = I_0 \rhd I_1 \rhd \cdots \rhd I_r = \{0\}$$
of $I$, in which every factor  is an elementary abelian $p$-group. Obviously,  each  $I_i/I_{i+1}$ is a simple $H$-module.  If $\sigma$ is trivial, then, by Corollary~\ref{inducedaction}, $\Lambda_H$ acts on $I$ via $\phi^{(\nu, \mu)}$. It is now not difficult to prove that  each $I_i / I_{i+1}$ is a simple $\Lambda_H$-module via the action $\phi^{(\nu, \mu)}$. Indeed, 	if  $K:=I_i / I_{i+1}$ is not simple $\Lambda_H$-module under $\phi^{(\nu, \mu)}$, then there exists a nonzero proper subgroup \(M < K\) which is invariant under \(\phi^{(\nu, \mu)}\).  
In particular,
	\[
	\phi^{(\nu, \mu)}_{(h,0)}(M)=\mu^{-1}_h(M)\subseteq M
	\quad\text{and}\quad
	\phi^{(\nu, \mu)}_{(0,h)}(M)=\nu_h(M)\subseteq M,
	\]
	which contradicts the hypothesis that $K$ is a simple $H$-module. The proof is complete.
\end{proof}

We remark that if a skew brace $H$ acts on an abelian group $I$ via $(\nu, \mu, \Id)$ such that $I$ is a simple $\Lambda_H$-module via the action $\phi^{(\nu, \mu)}$, then $I$ is a simple $H$-module. 

\begin{prop}
Let $E$ be a skew brace admitting a trivial brace $I$. Suppose that $I$ has a series
$$I = I_0 \trianglerighteq I_1 \trianglerighteq \cdots \trianglerighteq I_r = \{0\}$$
of its subgroups such that each $I_i$ is an ideal of $E$ for $i = 0,1 \ldots, r$. If there exist sub-skew braces $B_1, \ldots, B_r$ of $E$ such that $E = B_1+ I$, $B_1 \cap I = I_1$, and for each $i = 1, \ldots, r-1$, $B_i = B_{i+1} + I_i$ and $B_{i+1} \cap I_i = I_{i+1}$, then  $B_r$ is a complement of $I$ in $E$. 
\end{prop}

\begin{proof}
It is given that $E = B_1 + I$ and $B_i = B_{i+1} + I_i$ for each $i = 1, \ldots, r-1$. We obtain
	\[
	E = B_1 + I = B_2 + I_1 + I = B_2 +  I = \cdots = B_r + I.
	\]
We claim that  $B_r \cap I = \{0\}$. Contrarily, assume that $0 \ne  x \in B_r \cap I$. Then $x \in B_i$ for each  $i = 1, \ldots, r-1$, since $B_r \subseteq B_{r-1} \subseteq \cdots \subseteq B_1$. Therefore, we get
	\[
	x \in B_1 \cap I = I_1 \Rightarrow x \in B_2 \cap I_1 = I_2 \Rightarrow \cdots \Rightarrow x \in B_{r-1} \cap I_{r-2} = I_{r-1}.
	\]
Then $x \in B_r \cap I_{r-1} = \{0\}$, which contradicts the assumption that $x \ne 0$. Hence the claim follows, and the proof is complete.
\end{proof}





\noindent {\bf Acknowledgement:} The first-named author is thankful to ISI Delhi for providing a good research ambience. He thanks HRI Prayagraj for hosting him for two weeks in September 2025, during which the work got impetus via personal interaction with the other author. He also thanks Arpan Karnar for  fruitful discussions. The second-named author acknowledges the partial support of DST-SERB Grant MTR/2021/000285. 



\begin{thebibliography}{99}	
	
	\bibitem{DB18} D.~Bachiller, \textit{Extensions, matched products, and simple braces}, J. Pure Appl. Algebra \textbf{222} (2018), 1670--1691.
	
	\bibitem{Bach15} D.~Bachiller, \textit{Classification of braces of order $p^3$}, J. Pure Appl. Algebra \textbf{219}(8) (2015), 3568--3603.
	
%	\bibitem{BR21} A.~Ballester-Bolinches and R.~Esteban-Romero, \textit{Triply factorised groups and the structure of skew left braces}, Commun. Math. Stat. \textbf{10} (2022), 353--370.

	
	\bibitem{BEJP24}
	A.~Ballester-Bolinches,  R.~Esteban-Romero, P.~Jimenez-Seral and V.~Perez-Calabuig, 
	\textit{ Soluble skew left braces and soluble solutions of the Yang-Baxter equation}, 
Adv. Math. \textbf{455} (2024), Paper No. 109880, 27 pp.

\bibitem{BEFPT24}
A.~Ballester-Bolinches, A, R.~Esteban-Romero, M.~Ferrara, V.~Perez-Calabuig and M.~Trombetti,
\textit{Finite skew braces of square-free order and supersolubility}, 
Forum Math. Sigma 1\textbf{2} (2024), Paper No. e39, 33 pp.

\bibitem{BEP24}
A.~Ballester-Bolinches, A, R.~Esteban-Romero and V.~Perez-Calabuig,
 \emph{A Jordan-Holder theorem for skew left braces and their applications to multipermutation solutions of the Yang-Baxter equation},
Proc. Roy. Soc. Edinburgh Sect. A {\bf 154} (2024), 793-809.

%\bibitem{BNY20}  V.~G.~Bardakov, M.~V.~Neshchadim, and M.~K.~Yadav,  \textit{Computing skew left braces of small orders}, International Journal of Algebra and Computation \textbf{30}(4) (2020), 839--851.
	
	\bibitem{BS21} V.~G.~Bardakov and M.~Singh, \textit{ A Wells type exact sequence for non-degenerate unitary solutions of the Yang-Baxter equation}, Homology Homotopy Appl. \textbf{24} (2022), no. 2, 31--51.
	
	\bibitem{CSV18} F.~Cedó, A.~Smoktunowicz, and L.~Vendramin, \textit{Skew left braces of nilpotent type}, Proc. Lond. Math. Soc. \textbf{118} (2018), 1367--1392.
	
\bibitem{CFT25}
I. Collazo, M. Ferrara and M. Trombetti,  \textit{On derived-indecomposable solutions of the Yang-Baxter equation
Colazzo},  Publ. Mat. \textbf{69} (2025), 171--193.
	
	\bibitem{CDMFT25} 
	A.~Caranti, I.~Del~Corso, M.~Di~Matteo, M.~Ferrara, and M.~Trombetti, 
	\textit{A Sylow theorem for finite supersoluble skew braces}, 
	arXiv:2506.00940 (2025).
	
\bibitem{Evens72} L.~Evens, \textit{The Schur Multiplier of a Semi-Direct Product}, Illinois J. Math. \textbf{16}(1) (1972), 166--181.	

%\bibitem{TM98}	T.~Gateva-Ivanova, M.~ Van den Bergh, \textit{Semigroups of  I-type}, J. Algebra \textbf{206} (1998),97--112.

 \bibitem{GAP}
The GAP Group, \emph{Groups Algorithms and Programming,}  version 4.12.2 (2018), available at  http://www.gap-system.org.
	
\bibitem{GV17} L.~Guarnieri and L.~Vendramin, \textit{Skew braces and the Yang--Baxter equation}, Math. Comp. \textbf{86} (2017), no.~307, 2519--2534.
	
%\bibitem{MH59} M.~Hall, Jr., \textit{The Theory of Groups}, Macmillan, New York, 1959.
	
\bibitem{Hill72}
V.~E.~Hill, 
\emph{Split and minimal abelian extensions of finite groups},
Trans. Amer. Math. Soc. \textbf{172} (1972), 329--337.	
	
%\bibitem{JKA19}
%E. Jespers, L. Kubat and  A. Van Antwerpen, \textit{The structure monoid and algebra of a non-degenerate settheoretic solution of the Yang–Baxter equation}, Trans. Am. Math. Soc. \textbf{372} (10) (2019), 7191--7223
	
	
\bibitem{GK87}	G.~ Karpilovsky, \textit{The Schur multiplier.} London Math. Soc. Monogr. (N.S.) \textbf{2} (1987).

%\bibitem{KT24} Y.~Kozakai and C.~Tsang, \textit{Representation theory of skew braces}, Int. J. Group Theory \textbf{14}(3) (2025), 149--164.
	
	\bibitem{LV16} V.~Labed and L.~Vendramin, \textit{Cohomology and extensions of braces},  Pacific J. Math. \textbf{284} (2016), no. 1, 191--212.
	
%	\bibitem{LV23} T.~Letourmy and L.~Vendramin, \textit{Isoclinism of skew braces}, Bull. Lond. Math. Soc. \textbf{55}(6) (2023), 2891--2906.
	
	\bibitem{LV24} T.~Letourmy and L.~Vendramin, \textit{Schur covers of skew braces}, J. Algebra \textbf{644} (2024), 609--654.
	
	\bibitem{PSY18}
	I.B.S.~Passi, M.~ Singh, and M.K.~ Yadav, \textit{Automorphisms of finite groups}. Springer Monographs in Mathematics. Springer, Singapore, 2018. xix+217 pp.
	
	\bibitem{NR24} N.~Rathee, \textit{Extensions and Well's type exact sequence of skew braces}, J. Algebra Appl. \textbf{23}(1) (2024), 245009.
	
	\bibitem{RSU24} N.~Rathee, M.~Singh, and A.~Udeep, \textit{Representations of skew braces}, \href{http://arxiv.org/abs/2408.03766}{arXiv:2408.03766} (2024).
	
	\bibitem{RY24} N.~Rathee, M.K.~Yadav, \textit{Cohomology, Extensions and Automorphisms of Skew Braces}, J. Pure and Applied Algebra \textbf{228}(2) (2024), 107462.
	
%	\bibitem{WR07}
%	W. Rump, \textit{ Braces, radical rings, and the quantum Yang-Baxter equation}, J. Algebra \textbf{307} (2007), 153--170.
	
%\bibitem{KT72} K.~Tahara, \textit{On the second cohomology groups of semidirect products}, Math. Z. \textbf{129} (1972), 365--379.
	
\bibitem{MT23} Marco Trombetti, \textit{The structure skew brace associated with a finite non-degenerate solution of the Yang-Baxter equation is finitely presented}, Proc. AMS {\bf 152},  (2024), 573--583.
		
\bibitem{VK24}
L. Vendramin and  O. Konovalov, {\it Yang-Baxter}, version 0.10.6  (2024), available at https://www.gap-system.org/Packages/yangbaxter.html.

	
\end{thebibliography}



\end{document}






\begin{proof}
Let the  order  of $E$ be  \(p_1^{k_1} \cdots p_r^{k_r}\), where $p_i$'s are distinct primes and $k_i$'s and $r$ are natural numbers.  We proceed by the method of contradiction. So assume that $E$ is  the skew brace of the smallest order not admitting a sub-skew brace of order $p_j^{k_j}$ for at least one $j$, but all skew braces $A$ of order smaller then the order of $E$ admits Sylow $p$-sub-skew brace for every  prime $p$ dividing the order of $A$.  By Theorem \ref{minimumideal}, $E$ admits a minimal ideal $N$ of prime power order $p^k$ for a prime $p$ and a positive integer $k \ge 1$. So $p = p_i$ for some $i$. Consider the quotient skew brace $E/N$.  Since \(|E/N| < |E|\) and the hypothesis of the theorem holds true for $E/N$,  the assertion holds true for $E/N$.  

Let $\bar{K_j}$ denote a sub-skew brace  of $E/N$ having order $p_j^{k_j}$ for $j \ne i$. Note that, for the prime $p_i$, the sub-skew brace $K_i$ of $E/N$ has order $|\bar{K_i}| = p_i^{k_i-k}$. If $k_i = k$, then $N$ itself is the Sylow $p_i$-sub-skew brace of $E$ and we get an exact sequence $0 \to N \to E \to E/N \to 0$. Since $|N|$ is coprime to $|E/N|$, by Theorem \ref{main7}(3), $N$ admits a complement $L$ in $E$. Since $L < |E|$ and also the given hypothesis holds true for $L$, it admits Sylow $p_j$-sub-skew brace of $E$ for all $j \ne i$. Now assume that $k < k_i$.  Let $K_j$ be the pre-image of \( \bar{K_j} \), $1 \le j \le r$, in \( E \). Then, obviously, $K_i$ is a Sylow $p_i$-sub-skew brace of $E$. For $j \ne i$, we get the exact sequence $0 \to N \to K_j \to \bar{K_j} \to 0$. Since the order of $N$ is coprime to the order of 
$\bar{K_j}$, by Theorem \ref{main7}(3), $N$ admits a complement $L_j$ in $K_j$. It is now clear that $L_j$ is a Sylow $p_j$-sub-skew brace of $E$ for all $j \ne i$. Hence $E$ admits a 	Sylow $p$-sub-skew brace for every  prime divisor $p$ of $|E|$, which contradicts the minimality of $E$. The proof is complete.
\end{proof}
